% concatenated from main.tex
\documentclass[11pt,a4paper]{article}

\usepackage{flafter,amsmath,amssymb,latexsym,psfrag,graphicx,color,indentfirst,bm,amsthm}
\usepackage{mathrsfs}
\usepackage{bbm}
\usepackage{hyperref}
\usepackage{graphicx}
\usepackage{fullpage}
\usepackage{verbatim}
\usepackage[T1]{fontenc}
\usepackage[utf8]{inputenc}
\numberwithin{equation}{section}
\usepackage[margin=0.9in]{geometry}

% Needed by \ifthenelse in \todoautc
\usepackage{ifthen}
% Needed by the \begin{comment}...\end{comment} block in the bibliography
\usepackage{comment}

%margin/todo notes
\usepackage{todonotes}
\setlength{\marginparwidth}{2cm}
\newcommand{\todoautc}[3][]{%
	\ifthenelse{\equal{#1}{}}{\todo[size=\scriptsize]{{\bf#2} #3}{}}{\todo[backgroundcolor=#1,tickmarkheight=.5em,linecolor=orange,size=\scriptsize]{{\bf#2} #3}{}}%
}

\newcommand{\noteC}[1]{\todoautc[cyan!50]{}{#1}}
\newcommand{\noteT}[1]{\todoautc[none]{}{#1}}

% Theorem environments
\newtheorem{theorem}{Theorem}[section]
\newtheorem{lemma}[theorem]{Lemma}
\newtheorem{proposition}[theorem]{Proposition}
\newtheorem{corollary}[theorem]{Corollary}
\theoremstyle{definition}
\newtheorem{definition}[theorem]{Definition}
\newtheorem{remark}[theorem]{Remark}
\newtheorem{assumption}[theorem]{Assumption}

% Short-hands
\newcommand{\R}{\mathbb{R}}
\newcommand{\N}{\mathbb{N}}
\newcommand{\C}{\mathbb{C}}
\newcommand{\Om}{\Omega}
\newcommand{\eps}{\varepsilon}
\newcommand{\1}{\mathbbm{1}}
\newcommand{\dd}{\,\mathrm{d}}
\newcommand{\norm}[1]{\left\| #1 \right\|}
\newcommand{\ip}[2]{\left\langle #1, #2 \right\rangle}
\newcommand{\kappam}{\kappa_m}

\begin{document}

\title{%\LARGE \textbf{Feedback Control for Heat Tracking}\\
%\LARGE \textbf{Using Thermo-Plasmonic Nanoparticles}
\Large \textbf{Feedback Stabilization and Tracking for Heat Equations}\\ \textbf{Using Thermo-Plasmonic Nanoparticles as Actuators}
}

\author{Arpan Mukherjee\footnote{MSU-BIT SMBU Joint Research Center of Applied Mathematics, Shenzhen MSU-BIT University, Shenzhen, People's Republic of China (arpan.mukherjee@smbu.edu.cn and arpanmath99@alumni.iitm.ac.in).} ,\ Sérgio S. Rodrigues\thanks{  Center for Mathematics and Applications (NOVA Math) and Department of Mathematics, NOVA School of Science and Technology (NOVA FCT), NOVA University of Lisbon, Campus de Caparica
2829-516,  Portugal (ssi.rodrigues@fct.unl.pt).}
\ and
Mourad Sini\thanks{Radon institute, RICAM, Austrian Academy of Sciences, Altenbergerstrasse 69, 4040 Linz, Austria (mourad.sini@oeaw.ac.at).}}

\date{\today}

\maketitle

\begin{abstract}

We propose and analyze a strategy for tracking prescribed heat profiles in a
medium by using plasmonic nanoparticles as actuators. The starting point is a
thermo--plasmonic model coupling Maxwell's equations to a heat equation in a
homogeneous background containing a discrete collection of highly contrasting
nanoparticles. A recently derived time-domain discrete effective model in \cite{CaoMukherjeeSini2025} shows
that, outside the particle region, the heat generated by the nanoparticles is
well approximated by a superposition of heat kernels centered at the particle
locations, with time-dependent amplitudes solving a coupled Volterra system.

\medskip
\noindent
Our first step is to recast this discrete thermo--plasmonic description as a
heat equation with finitely many point actuators on a bounded control domain
$\Omega$. We then design and analyze a heat-profile
tracking feedback law based on pointwise evaluations of $\mathcal A^{-1}y$,
where $\mathcal A=I-A_0$ and $A_0$ is the Neumann diffusion operator. The
closed-loop dynamics is treated in the natural $V'$ framework, where $V := D(\mathcal A)$, and we obtain
exponential stabilization of the tracking error in $V'$ by relying on the
distribution-actuator theory in \cite{KunRodWal24-cocv}. For general
(non-equilibrium) reference profiles we incorporate a constant feedforward and
a fixed-point pre-compensation on a low-mode space $X_N$, which yields exact
matching of the steady target on $X_N$ together with an explicit estimate of
the remaining tail mismatch.

\medskip
\noindent
Finally, we couple the thermo--plasmonic effective model with the feedback
design and quantify the discrepancy between the ideal closed-loop temperature
and the physically realized one. The resulting tracking bound separates (i) a
projection error due to restricting the abstract feedback input to a
realizable finite-dimensional space and (ii) a thermo--plasmonic remainder
that vanishes as the particle size $\delta\to0$. We also make explicit how
plasmonic resonances enter the effective actuation map and how resonant tuning
can improve control authority and conditioning. This provides a rigorous
framework for using plasmonic nanoparticles as actuators for heat-profile
tracking, with potential relevance for nano-localized photothermal therapy and
temperature management in nano-engineered materials.



\end{abstract}

\section{Introduction}

\subsection{Thermo-plasmonic heating and effective medium theory}

\noindent Metallic (plasmonic) nanoparticles exhibit strong absorption and local field
enhancement when illuminated near their plasmonic resonances. This leads to
highly localized heat generation that can be exploited in a wide range of
applications, including photothermal cancer therapy, thermal imaging,
nano-manipulation and temperature management in micro- and nano-devices; see,
for instance, \cite{BaffouQuidant2013,BaffouNM2020,Chuang2022,
Jorgensen2016,Kim2023, Vines2019}. We discuss in Subsection~\ref{subsec:plasmonic-resonances} how resonant tuning enters the asymptotic modeling and improves the effective actuation strength.

\noindent From a modeling viewpoint, quantitative descriptions of light-to-heat
conversion in plasmonic systems typically couple a time-harmonic Maxwell
problem, governing the electromagnetic field, to a heat equation with
volumetric source term representing ohmic losses. In the physical and
engineering literature a variety of models have been proposed to describe the
temperature rise around single particles and suspensions of particles, often
relying on phenomenological or numerical approaches
\cite{BaffouQuidant2013,BaffouNM2020,PhotothermalJAP2016,WuSheldon2023}. On
the mathematical side, asymptotic and effective-medium descriptions have been
developed for plasmonic nanoparticles in the quasi-static and subwavelength
regimes, with a particular emphasis on the role of resonances and high
contrast; see, e.g.,
\cite{AmmariMillienRuizZhang2017,AmmariRomeroRuiz2018,
CaoMukherjeeSini2025,Mukherjee2022} and the references therein.

\noindent In \cite{CaoMukherjeeSini2025}, a rigorous time-domain effective model for
heat generation by a finite collection of plasmonic nanoparticles in an
otherwise homogeneous background was obtained, under subwavelength and
high-contrast scaling assumptions. The main result shows that the temperature
difference $w$ between the actual temperature and the background solution
without particles can be approximated, outside the particle region, by a
finite superposition of heat kernels centered at the nanoparticle centers,
with time-dependent amplitudes $\sigma_i$ solving a Volterra-type integral
system. The error is quantified in suitable norms, and the amplitudes are
shown to depend linearly (up to higher-order terms) on the internal
electromagnetic energy generated by the incident illumination. This discrete
description is based on a Foldy--Lax point-approximation framework applied to
the underlying Maxwell scattering problem.

\noindent The representation of $w$ as a superposition of heat kernels localized at the
particle centers is the key structural feature that we exploit here. It
suggests that, at least in the effective description, a suitably designed
ensemble of plasmonic nanoparticles can be interpreted as a system of point
actuators for the heat equation, remotely driven by electromagnetic inputs.

\subsection{Heat-profile tracking and point actuators}

\noindent  Independently of the thermo--plasmonic context, the control and stabilization
of linear parabolic equations has been extensively studied; see, e.g.,
\cite{Barbu2018,KrsticSmyshlyaev2008,Pazy1983,TucsnakWeiss2009,RaymondBook}
and the references therein. A particularly relevant class of problems for the
present work involves finite-dimensional actuation and observation, where a
heat equation is driven by a finite number of internal or boundary actuators
and observed through finitely many output functionals.

\noindent Pointwise (Dirac-type) actuation and observation models are natural
idealizations when the actual sources and sensors are localized on spatial
scales much smaller than the size of the domain. In the context of the heat
equation, controllability and stabilization issues for point actuators and for
controls acting on small or moving subsets have been investigated in various
settings; see, for instance, \cite{CasZuazua,
FardigolaKhalina2022,Hamidoglu2013,Letrouit2019, Martinez2015,Rodrigues2021} and the references
therein. The analysis of such systems typically relies on semigroup methods,
spectral decompositions and Carleman estimates.
\bigskip

\noindent In the present paper, we consider a controlled heat equation on a bounded
domain $\Omega\subset\R^d$, mainly $d=3$, with homogeneous Neumann boundary conditions, driven by finitely many point actuators. We design an output feedback based on
\emph{elliptically regularized} point measurements, namely pointwise
evaluations of $\mathcal A^{-1}y$, where $\mathcal A=I-A_0$ and $A_0$ is the
Neumann diffusion operator. For a prescribed reference profile $y_r$, the
feedback has the form
\[
u_j(t) = (u_r)_j-\lambda\Big[(\mathcal A^{-1}y)(x_j,t) - (\mathcal A^{-1}y_r)(x_j)\Big],
\qquad j=1,\dots,M,
\]
where $\lambda>0$ is a gain and $u_r\in\R^M$ is a constant feedforward used to
compensate the fact that non-constant references are not equilibria of the
free Neumann heat dynamics (and which is further refined via a fixed-point
pre-compensation on low modes; see Section~\ref{subsec:fixed-point-bias}).
Working in the $V'$ framework, with $V:=D(\mathcal A)$, we use the stabilization theory in \cite{KunRodWal24-cocv} to obtain exponential convergence of the tracking
error in $V'$ for this class of feedbacks; see Section~\ref{sec:stab-KRW}.



\subsection{Objective and main contributions}
\noindent
The goal of this work is to connect rigorously the discrete thermo--plasmonic
effective model of \cite{CaoMukherjeeSini2025} with heat-profile tracking
feedback laws for a heat equation with point actuators of \cite{KunRodWal24-cocv} to propose and analyse a strategy for tracking prescribed heat profiles in a medium by using plasmonic nanoparticles as actuators. We achieve this goal with the following steps. 
\begin{itemize}
\item[(i)] We first recall the discrete thermo--plasmonic
effective model of \cite{CaoMukherjeeSini2025}, in Section \ref{sec:model}, and then recast it  as a heat equation
on a bounded domain driven by finitely many point actuators, in a
functional-analytic setting suitable for control design. In particular, we
identify the effective point-source amplitudes in terms of the Volterra system
for the nanoparticle amplitudes and provide a quantitative justification for
restricting the whole-space heat propagation in $\R^3$ to a bounded Neumann
control domain on compact subsets away from $\partial\Omega$ (Lemma~\ref{lem:R3-to-Omega-restriction}). Precisely, the temperature field satisfies in $\R^3$
\[
\partial_t y-\kappa_m\Delta y = Q_\delta(p),
\]
where $Q_\delta(p)$ is the nanoparticle-induced volumetric source supported in $\delta$-type sources located of the particle centers. Upon restriction to a bounded Neumann control domain $\Omega$ containing the actuator locations $\{x_j\}_{j=1}^M$, we recast the dynamics as
\[
\dot y = A_0 y + Bu,
\qquad
Bu=\sum_{j=1}^M u_j\,\delta_{x_j},
\]
posed in the extrapolation space $V'$, where $V: =D(\mathcal A)$, with $\mathcal A=I-A_0$. Then, in Section \ref{sec:point-actuators}, we establish well-posedness of this problem improving the related results in \cite{KunRodWal24-cocv}, allowing more general shapes $\Omega$ with $C^{1,1}$-regularity; see Proposition \ref{prop:wp-Dirac}.




\item[(ii)] We state, in Section \ref{sec:feedback}, the feedback design and stability analysis by projecting into finite-dimensional spaces with a control of the tails. Then, in Section \ref{sec:stab-KRW}, we establish well-posedness and basic stability properties of a
feedback law for heat-profile tracking based on pointwise evaluations of
$\mathcal A^{-1}y$, interpreted as a closed-loop parabolic system with point
actuators. We place the closed loop in the distribution-actuator framework of
\cite{KunRodWal24-cocv} in the state space $V'$ and obtain exponential
convergence to an equilibrium. For non-equilibrium target profiles we
incorporate a feedforward term and a fixed-point pre-compensation on a low-mode
space $X_N$, yielding exact matching of the first $N$ Neumann modes of the
steady reached profile and an explicit tail estimate; see
Section~\ref{subsec:fixed-point-bias}. Namely, the feedback uses the elliptically regularized point-evaluation operator
\[
C:V'\to\R^M,
\qquad
Cz=\big((\mathcal A^{-1}z)(x_j)\big)_{j=1}^M,
\]
and, for references $y_r\in X_N:=\mathrm{span}\{\varphi_1,\ldots,\varphi_N\}$, takes the form
\[
u(t)=u_r-\lambda\,C\big(y(t)-y_r\big),
\]
with $u_r\in\R^M$ chosen so that $P_N(A_0 y_r+Bu_r)=0$.


\item[(iii)] We couple the effective thermo--plasmonic model with the feedback
design and quantify the difference between the ideal closed-loop
temperature (in the abstract control model) and the temperature realized
physically through plasmonic actuation, see Section \ref{sec:coupling}. At this level, we obtain a tracking
bound that separates a projection error (restriction of the abstract feedback
to a realizable finite-dimensional control space) from a thermo--plasmonic
modeling remainder that vanishes as $\delta\to0$, where $\delta$ is the maximum radius of the nanoparticles, and we explain how plasmonic
resonances enter the actuation map and can be exploited to improve control
authority and conditioning (Subsection~\ref{subsec:plasmonic-resonances}). Precisely, at the level of the thermo-plasmonic effective model, the nanoparticle-induced heat input admits the representation
\[
G_\delta(p)=Kp+R_\delta p,
\qquad
\|R_\delta\|_{\mathcal L(U,L^2(0,T;\R^M))}\xrightarrow{\,\delta\to0\,}0,
\]
where $K\in\mathcal L\bigl(U,L^2(0,T;\R^M)\bigr)$ is a bounded linear actuation operator mapping illumination controls to effective point-source amplitudes.


\item[(iv)] In Section~\ref{sec:robust-Vprime}, we derive robustness bounds that quantify the effect of implementing the feedback only approximately. More precisely, if the realized heat input differs from the ideal Dirac-actuator feedback by a discrepancy term in $D(\mathcal A)'$, then the induced tracking deviation is controlled in $C([0,T];V')\cap L^2(0,T;H)$ on finite horizons, with an additional $L^2(\Omega)$ estimate on $[t_0,T]$ for any $t_0>0$.

\item[(v)] In Section~\ref{rem:link-original-problem}, we explain how the main tracking theorem translates back to the original thermo--plasmonic control problem. The resulting estimate separates (a) the abstract closed-loop tracking error, (b) a projection error due to restricting the feedback input to a realizable subspace $Y_0$, and (c) a thermo--plasmonic realization remainder that vanishes as $\delta\to0$.
\par\smallskip
In particular, the implemented tracking error admits a decomposition of the form
\[
y-y_r=\mathcal E_{\mathrm{proj}}+\mathcal E_{\delta},
\]
where $\mathcal E_{\mathrm{proj}}$ corresponds to restricting the ideal feedback input to a realizable subspace $Y_0$, and $\mathcal E_{\delta}$ is induced by the thermo--plasmonic remainder.


\item[(vi)] In Section~\ref{sec:algorithm}, we provide an algorithmic design that turns the theory into a computable feedback loop: we build an implementable feedback law (including the feedforward and fixed-point correction), identify/assemble a finite-dimensional representation of the thermo--plasmonic map $\mathcal K$ on $U_0\to Y_0$, and use a (pseudo-)inverse to recover admissible incident-field intensities that realize the desired reduced heat input.

\item[(vii)] Finally, in Section~\ref{appendix},  as an appendix, we collect the technical ingredients used throughout the analysis, including verifications of the standing spectral and actuator assumptions (on $C^{1,1}$ domains) and sufficient conditions ensuring solvability of the fixed-point correction for general reference profiles.

\end{itemize}





\noindent The main message is that---under the scaling and smallness assumptions
guaranteeing the validity of the discrete thermo--plasmonic model, and under
standard controllability-type conditions on the actuator locations and on the illumination patterns---it is meaningful to interpret a configuration of
plasmonic nanoparticles as a system of actuators for the heat equation. By exciting the nanoparticles with suitably designed electromagnetic incident fields, one
can approximate abstract feedback laws and hence track a desired heat profile. The strategy yields exact matching of the desired profile in a prescribed finite-dimensional spectral subspace (by projection) and an explicit estimate of the remaining tail mismatch; see Section~\ref{subsec:fixed-point-bias}.
\bigskip



\noindent We finish this introduction by mentioning two idealizations that are used in the feedback design. First, the coefficient vector $p(t):=(p_l(t))^m_{l=1}\in\mathbb{R}^m$ is treated as a signed control acting linearly through the illumination dictionary (coming from the electromagnetic fields generated by the injected nanoparticles), but physically, throught the electromagnetism model, it appears of the form $p_\ell(t)=|a_\ell(t)|^2\ge 0$ (where $a_l$ are related to the electric field intensity).
This can be handled by restricting to effective inputs attainable with $p\in (\mathbb{R}_+)^m$, see Remark \ref{Positivity-p} for more details.
Second, the output map is taken as the elliptically regularized point evaluation $Cz=\big((A^{-1}z)(x_j)\big)_{j=1}^M$, which is well-defined on the natural state space for Dirac-actuated heat dynamics and enables stability estimates in $V'$ (and hence in $L^2$ for $t\ge t_0>0$).
However, in practice, this corresponds to one more step of reconstructing these regularized outputs from physically available measurements (as those collected as local averages or only near the boundary of $\Omega$, for instance). These two issues, i.e. incorporating explicit nonnegativity and realistic sensing, will be discussed in a forthcoming work.

 

\section{Thermo-plasmonic model and discrete effective description}
\label{sec:model}
\noindent
We briefly recall the thermo--plasmonic model and the discrete effective
description from \cite{CaoMukherjeeSini2025}. We only state the parts needed
for the control analysis; the reader is referred to
\cite{CaoMukherjeeSini2025} for details and precise assumptions.

\subsection{Coupled Maxwell--heat model in $\R^3$}
\noindent
The physical space is $\R^3$ and the medium consists of a homogeneous
background characterized by constant volumetric heat capacity $c_m>0$,
thermal conductivity $\gamma_m>0$, and permittivity $\eps_m>0$, together with
$M$ plasmonic nanoparticles
\[
D_i := z_i + \delta B_i,\qquad i=1,\dots,M,
\]
where $z_i\in\R^3$ are the centers, $\delta\ll1$ is a small scaling
parameter, and the shapes $B_i$ are bounded $C^2$ domains with
$\mathrm{vol}(B_i)\sim1$. The particle material is characterized by
$c_p,\gamma_p,\eps_p$ and we denote the piecewise-constant coefficients
\[
c_v := c_p\chi_D + c_m\chi_{\R^3\setminus D},\quad
\gamma := \gamma_p\chi_D + \gamma_m\chi_{\R^3\setminus D},\quad
\eps := \eps_p\chi_D + \eps_m\chi_{\R^3\setminus D},
\]
where $D:=\bigcup_{i=1}^M D_i$.
\newline
The electric field $E$ solves a time-harmonic Maxwell scattering problem at
frequency $k$,
\begin{equation}
\label{eq:maxwell}
\begin{cases}
\operatorname{curl}\operatorname{curl}E - k^2\eps E = 0 & \text{in }\R^3,\\
E = E^{\mathrm{in}} + E^{\mathrm{sc}},\\
E^{\mathrm{sc}}\ \text{satisfies the Silver--M\"uller condition,}
\end{cases}
\end{equation}
for a given incident field $E^{\mathrm{in}}$. The choice of $\eps_p$
follows a Drude model and the frequency $k$ is chosen near a plasmonic
resonance; see \cite{AmmariRomeroRuiz2018,
CaoMukherjeeSini2025, Mukherjee2022}.
\newline
The temperature $u=u(x,t)$ is governed by a heat equation for $x\in\R^3$ and
$t>0$:
\begin{equation}
\label{eq:heat-transmission}
\begin{cases}
c_v(x)\,\partial_t u - \nabla\cdot(\gamma(x)\nabla u) = J(x,t)
& \text{in }\R^3\times(0,T),\\[0.3em]
u(x,0) = 0 & \text{in }\R^3,
\end{cases}
\end{equation}
where the source term represents volumetric heating due to ohmic losses,
\begin{equation}\label{eq:heat=J}
J(x,t) := \frac{k}{2\pi}\Im\eps(x)\,|E(x)|^2 f(t),
\end{equation}
with a causal temporal profile $f\in H^{r+1/2}_{0,\sigma}(0,T)$ and where~$\Im \eps$ denotes the imaginary part of the permittivity $\eps$, which is taken positive valued and supported on the nanoparticles, meaning that the background medium (outside the cluster of nanoparticles) is lossless; see
\cite{CaoMukherjeeSini2025} for more details and the precise functional framework. The real valued 
constants
\[
\kappam := \frac{c_m}{\gamma_m},\qquad \kappa_p := \frac{c_p}{\gamma_p}
\]
denote the diffusion constants in the background and inside the particles.
\newline
The main object of interest is the temperature difference
$w:=u-u^{(0)}$, where $u^{(0)}$ is the solution of the homogeneous-background
problem (i.e.\ without particles). This difference satisfies a parabolic
transmission problem with source terms supported in the nanoparticle region
$D$.

\subsection{Discrete thermo-plasmonic effective model}
\noindent
We denote by $\Phi^{(m)}(x,t;y,\tau)$ the fundamental solution of the
homogeneous heat operator $\partial_t-\kappam\Delta$ in $\R^3$:
\begin{equation}
\label{eq:heat-kernel}
\Phi^{(m)}(x,t;y,\tau)
= \frac{1}{\big(4\pi\kappam(t-\tau)\big)^{3/2}}
\exp\left(-\frac{|x-y|^2}{4\kappam(t-\tau)}\right)
\chi_{(0,\infty)}(t-\tau).
\end{equation}
\noindent
Under subwavelength, separation and high-contrast assumptions, the main
effective result in \cite{CaoMukherjeeSini2025} can be summarized as follows.

\begin{theorem}[Discrete thermo-plasmonic effective model]
\label{thm:discrete-effective}
Assume the scaling and regularity hypotheses of \cite{CaoMukherjeeSini2025},
in particular the subwavelength, separation, and high-contrast conditions on
$(c_p,\gamma_p,\eps_p)$ and the nanoparticles $D_i = z_i + \delta B_i$.
Then there exist functions
\[
\sigma_i\in H^1(0,T),\qquad i=1,\dots,M,
\]
such that the temperature difference $w=u-u^{(0)}$ satisfies, for
$x\in\R^3\setminus D$ and $t\in(0,T)$,
\begin{equation}
\label{eq:discrete-expansion}
w(x,t) = -\sum_{i=1}^M \frac{\alpha_i}{c_m}\int_0^t
\Phi^{(m)}(x,t;z_i,\tau)\,\sigma_i(\tau)\,\dd\tau + \mathcal{E}(x,t),
\end{equation}
where $\alpha_i:=\gamma_{p,i}-\gamma_m$ denotes the heat-conductivity
contrast, and the error $\mathcal{E}$ admits the bound
\[
\norm{\mathcal{E}(\cdot,t)}_{L^2(K)} \le C M\delta^\mu,\qquad t\in[0,T],
\]
for each compact $K\subset\R^3\setminus D$, for some $\mu>0$ depending on
the scaling exponents and on $K$.
\newline
Moreover, the amplitudes $\sigma_i$ solve a Volterra integral system of the
form
\begin{equation}
\label{eq:Volterra-system}
\sigma_i(t) + \sum_{j\ne i}\beta_{ij} \int_0^t
\partial_t\Phi^{(m)}(z_i,t;z_j,\tau)\,\sigma_j(\tau)\,\dd\tau
= F_i(t) + r_i(t),
\end{equation}
where the interaction coefficients $\beta_{ij}$ depend on the geometry and
contrast parameters, the forcing terms $F_i$ depend linearly on the internal
electromagnetic energy $|E|^2$ in $D_i$, and the remainders $r_i$ satisfy
\[
\sum_{i=1}^M \norm{r_i}_{H^1(0,T)} \le C M\delta^{\mu'}
\]
for some $\mu'>0$. The system \eqref{eq:Volterra-system} is uniquely
solvable in $(H^1(0,T))^M$ and the $\sigma_i$ satisfy uniform a priori
bounds in terms of the $F_i$.
\end{theorem}

\begin{remark}
The precise exponents $\mu,\mu'$ and the constants in the above estimates
depend on the detailed asymptotic regime chosen in \cite{CaoMukherjeeSini2025}.
For the control-oriented discussion below, it is enough to retain that the
errors vanish as $\delta\to0$ and that the $\sigma_i$ are uniformly bounded
in $H^1(0,T)$ under the considered scaling regime.
\end{remark}


\subsection{Plasmonic resonances and actuation strength}
\label{subsec:plasmonic-resonances}
\noindent
The efficiency of the actuation mechanism is based on the magnitude of the internal
electromagnetic energy $\int_{D_i}|E|^2$ entering the forcing terms $F_i$ in
\eqref{eq:Volterra-system}. This quantity is strongly amplified when the
illumination frequency is tuned near a \emph{plasmonic resonance} of the nanoparticles.
We recall below how resonances enter the modeling in the language of the
Maxwell--Lippmann--Schwinger formulation, based on the magnetization operator and
the Helmholtz decomposition as in \cite{CaoMukherjeeSini2025}.

\medskip
\noindent
\textbf{Magnetization operator and resonant resolvent.}
In the time-harmonic Maxwell problem \eqref{eq:maxwell} with piecewise-constant
permittivity $\eps=\eps_p\chi_D+\eps_m\chi_{\R^3\setminus D}$, the electric field
satisfies a Maxwell--Lippmann--Schwinger equation of the form
\begin{equation}
\label{eq:LS-magnetization}
E(x) + \eta \sum_{j=1}^M M^{(k)}_{D_j}[E](x) - k^2\eta \sum_{j=1}^M N^{(k)}_{D_j}[E](x) = E^{\mathrm{in}}(x),
\qquad x\in D,
\end{equation}
where $\eta:=\eps_p-\eps_m$ and, with $G^{(k)}$ the Helmholtz Green's function,
\[
M^{(k)}_{D}[f](x):=\nabla \int_D \nabla G^{(k)}(x,y)\cdot f(y)\,dy,
\qquad
N^{(k)}_{D}[f](x):=\int_D G^{(k)}(x,y)\, f(y)\,dy.
\]
To analyze \eqref{eq:LS-magnetization} in each particle $D_i=\delta B_i+z_i$, we use
the Helmholtz decomposition
\begin{equation}
\label{eq:Helmholtz-decomp}
L^2(D_i)^3=H_0(\mathrm{div}\,0,D_i)\oplus H_0(\mathrm{curl}\,0,D_i)\oplus \nabla \mathrm{Harm}(D_i),
\end{equation}
and denote by $\mathbb P^{(3)}$ the orthogonal projector onto $\nabla\mathrm{Harm}(D_i)$.
In the subwavelength regime $\delta k\ll 1$, the dominant amplification occurs on
$\nabla\mathrm{Harm}(D_i)$, where the quasi-static magnetization operator
$M^{(0)}_{B_i}:\nabla\mathrm{Harm}(B_i)\to\nabla\mathrm{Harm}(B_i)$ admits an
orthonormal eigenbasis $\{e^{(3)}_{i,n}\}_{n\in\mathbb N}$ with eigenvalues
$\{\lambda^{(3)}_{i,n}\}_{n\in\mathbb N}$, so that
\[
M^{(0)}_{B_i}[e^{(3)}_{i,n}] = \lambda^{(3)}_{i,n}\, e^{(3)}_{i,n}.
\]
Projecting \eqref{eq:LS-magnetization} onto $\nabla\mathrm{Harm}(D_i)$ and scaling
$D_i$ to $B_i$ yields, at leading order,
\begin{equation}
\label{eq:mode-balance}
\bigl(I+\eta M^{(0)}_{B_i}\bigr)\,\mathbb P^{(3)}\widehat E_i
\;\approx\; \mathbb P^{(3)}\widehat E^{\mathrm{in}}_i \;+\; \text{(particle--particle interaction terms)}.
\end{equation}
Consequently, the modal coefficients satisfy
\[
\langle \mathbb P^{(3)}\widehat E_i, e^{(3)}_{i,n}\rangle
\;\approx\; \frac{1}{1+\eta\lambda^{(3)}_{i,n}}\,
\Big(\langle \mathbb P^{(3)}\widehat E^{\mathrm{in}}_i, e^{(3)}_{i,n}\rangle+\cdots\Big).
\]
Plasmonic resonances correspond to frequencies (through $\eps_p(k)$, hence $\eta$)
for which a \emph{dispersion relation} $1+\eta\lambda^{(3)}_{i,n}=0$ is (approximately)
satisfied for some index $n=n_0$. In the asymptotic regime used in
\cite{CaoMukherjeeSini2025}, the frequency is chosen so that
\begin{equation}
\label{eq:resonance-scaling}
|1+\eta\lambda^{(3)}_{i,n}|\sim
\begin{cases}
\delta^{h}, & n=n_0,\\
1, & n\neq n_0,
\end{cases}
\end{equation}
which implies the resolvent growth
$\|(I+\eta M^{(0)}_{B_i})^{-1}\|_{\mathcal L(L^2(B_i)^3)}\lesssim \delta^{-h}$ on
$\nabla\mathrm{Harm}(B_i)$ and yields the near-field enhancement
$\|\mathbb P^{(3)}E\|_{L^2(D_i)} \sim \delta^{3/2-h}$ (and similarly in $L^4$), up to
interaction corrections.

\medskip
\noindent
\textbf{Polarization tensor and absorbed power.}
The effective electromagnetic coupling entering the discrete heat model
\eqref{eq:Volterra-system} is encoded by the polarization matrix $P_{D_i}$ appearing
in the Foldy--Lax system for the projected moments of $E$ (see \cite{CaoMukherjeeSini2025}).
In the magnetization-operator spectral variables, one has the representation
\begin{equation}
\label{eq:polarization-magnetization}
P_{D_i} \;=\; \delta^3 \sum_{n\in\mathbb N}\frac{1}{1+\eta\lambda^{(3)}_{i,n}}
\left(\int_{B_i} e^{(3)}_{i,n}(\xi)\,d\xi\right)\otimes
\left(\int_{B_i} e^{(3)}_{i,n}(\xi)\,d\xi\right)
\;+\; O(\delta^3),
\end{equation}
so that, under \eqref{eq:resonance-scaling}, the dominant resonant contribution scales like
$P_{D_i}\sim \delta^{3-h} P_{B_i}$ for an explicitly computable shape-dependent matrix
$P_{B_i}$.
Since the ohmic source term in \eqref{eq:heat-transmission} is
$J(x,t)=\frac{k}{2\pi}\Im\eps(x)\,|E(x)|^2 f(t)$, the total absorbed power in $D_i$ satisfies
\begin{equation}\label{Q_i-s}
Q_i \;:=\; \frac{k}{2\pi}\Im\eps_p(k)\int_{D_i}|E(x)|^2\,dx
\;\approx\; \Im\eps_p(k)\, \delta^{3-2h}\,\mathfrak A_i \,\big|E^{\mathrm{in}}(z_i)\big|^2,
\end{equation}
with an amplification factor $\mathfrak A_i$ determined by the resonant mode(s) in
\eqref{eq:polarization-magnetization} and by the inter-particle Foldy--Lax coupling.
In particular, tuning the illumination near resonance increases $Q_i$ by orders of
magnitude compared to off-resonant operation.



\medskip
\noindent
\textbf{Consequences for control conditioning and tracking.}
In the control reduction developed below, the finite-dimensional map from
illumination parameters $p\in U_0$ to effective heat inputs $u\in Y_0$ is encoded by
the matrix/operator $K_0$ (Section~\ref{sec:coupling}). The resonant scaling
$P_{D_i}\sim \delta^{3-h}$ increases the gain of each actuator channel and typically
improves the conditioning of $K_0$ (larger effective singular values), which:
\begin{itemize}
\item reduces amplification in the right-inverse $K_0^\dagger$ needed to realize a
prescribed $u_{\mathrm{des}}\in Y_0$,
\item improves robustness with respect to modeling and measurement errors, and
\item allows one to enrich the realizable control space $Y_0$ (hence reduce projection
error) for a fixed admissible illumination energy.
\end{itemize}
Thus, plasmonic resonances are not only the physical origin of strong local heating,
but also a mathematical lever to strengthen actuation and improve tracking accuracy.


\section{Heat equation with point actuators on a bounded domain}
\label{sec:point-actuators}
\noindent
We now introduce the abstract control model for the heat equation with point
actuators, in a bounded domain $\Om\subset\R^3$. This section is purely
functional-analytic and independent of the thermo--plasmonic model, except
for the fact that the background diffusion coefficient is $\kappam$.

\subsection{Physical model and abstract reformulation}
\noindent
Let $\Om\subset\R^3$ be a bounded, connected, $C^2$ domain. The restriction
of the background heat equation to $\Om$ (with homogeneous Neumann boundary
conditions) reads
\begin{equation}
\label{eq:3.1}
\begin{cases}
\partial_t y - \kappam\Delta y
= \displaystyle\sum_{j=1}^M u_j(t)\,\delta_{x_j}
& \text{in }\Om\times(0,T),\\[0.3em]
\partial_\nu y = 0 & \text{on }\partial\Om\times(0,T),\\[0.3em]
y(\cdot,0) = y_0 & \text{in }\Om.
\end{cases}
\end{equation}

\begin{lemma}[Restriction from $\R^3$ to a bounded control domain]
\label{lem:R3-to-Omega-restriction}
Let $\Omega\subset\R^3$ be a bounded $C^2$ domain and let $B\Subset\Omega$ be an
open set satisfying $\mathrm{dist}(B,\partial\Omega)=:d>0$. Consider a source
term $f(\cdot,t)$ supported in $B$ for each $t\in(0,T)$ and let $w$ be the
whole-space solution of
\[
\partial_t w-\kappam\Delta w=f\quad\text{in }\R^3\times(0,T),\qquad w(\cdot,0)=0,
\]
while $y$ solves the Neumann problem
\[
\partial_t y-\kappam\Delta y=f\quad\text{in }\Omega\times(0,T),\qquad
\partial_\nu y=0\ \text{on }\partial\Omega\times(0,T),\qquad y(\cdot,0)=0.
\]
Then, for all $t\in(0,T]$,
\begin{equation}
\label{eq:R3-to-Omega-bound}
\|w(t)-y(t)\|_{L^2(B)}
\le C\int_0^t (t-\tau)^{-3/4}\exp\!\Big(-\frac{c\,d^2}{t-\tau}\Big)\,
\|f(\tau)\|_{L^2(B)}\,\dd\tau,
\end{equation}
where $C,c>0$ depend only on $\Omega$, $B$ and $\kappam$. In particular,
\begin{equation}
\label{eq:R3-to-Omega-simple}
\sup_{t\in[0,T]}\|w(t)-y(t)\|_{L^2(B)}
\le C\,\exp\!\Big(-\frac{c\,d^2}{T}\Big)\,\|f\|_{L^1(0,T;L^2(B))}.
\end{equation}
\end{lemma}

\begin{proof}[Proof sketch]
Both solutions admit Duhamel representations in terms of the whole-space heat
kernel $\Phi^{(m)}$ and the Neumann heat kernel $G_\Omega$ on $\Omega$.
For $x\in B$ one has
\[
(w-y)(x,t)=\int_0^t\!\int_{B}\big(\Phi^{(m)}(x,t;\xi,\tau)-G_\Omega(x,\xi,t-\tau)\big)\,
f(\xi,\tau)\,\dd\xi\,\dd\tau.
\]
Standard Gaussian bounds for $G_\Omega$ and off-diagonal estimates for the
difference $G_\Omega-\Phi^{(m)}$ when $x,\xi\in B$ yield
\[
|\Phi^{(m)}(x,t;\xi,\tau)-G_\Omega(x,\xi,t-\tau)|
\le C(t-\tau)^{-3/2}\exp\!\Big(-\frac{c\,d^2}{t-\tau}\Big),
\]
which implies \eqref{eq:R3-to-Omega-bound} by Young's inequality. Estimate
\eqref{eq:R3-to-Omega-simple} follows by bounding the convolution kernel on
$(0,T]$.
\end{proof}



\begin{remark}
Lemma~\ref{lem:R3-to-Omega-restriction} provides a quantitative justification
for replacing the whole-space heat propagation by a bounded-domain Neumann
model when actuators, nanoparticles and the region of interest are contained
in a subdomain $B$ at positive distance from the boundary. In particular, for
time horizons $T\ll d^2$ the restriction error is exponentially small in
$d^2/T$.
\end{remark}
%--------------------------------------------------------
\subsection{Dirac actuators and low-regularity well-posedness}
\label{subsec:wp-Dirac-KRW}
%--------------------------------------------------------
\noindent
We refer to \cite{KunRodWal24-cocv} for distribution-type actuator frameworks
in closely related parabolic stabilization settings.
\newline
We set
\begin{equation}
\label{eq:3.2}
H := L^2(\Om),\qquad
A_0 y := \kappam\Delta y,\quad
D(A_0) := \{ y\in H^2(\Om): \partial_\nu y = 0 \text{ on }\partial\Om\}.
\end{equation}
Then $A_0$ generates an analytic contraction semigroup
$S(t):H\to H$ for $t\ge0$.
\newline
The control operator is defined by
\begin{equation}
\label{eq:3.3}
B:\R^M\to D(A_0)',\qquad
B u := \sum_{j=1}^M u_j\,\delta_{x_j}.
\end{equation}
With this notation, \eqref{eq:3.1} is equivalent to the abstract system
\begin{equation}
\label{eq:3.4}
\dot y(t) = A_0 y(t) + B u(t),\qquad y(0)=y_0.
\end{equation}
Let $H_{-1}$ be the extrapolation space associated with $A_0$ (completion of $H$ with respect to
$\|x\|_{H_{-1}} := \|(I-A_0)^{-1} x\|_{H}$).
\newline
Introduce the coercive elliptic operator
\[
\mathcal A := I - A_0 = I - \kappa_m \Delta,
\qquad D(\mathcal A)=D(A_0),
\]
which is an isomorphism $\mathcal A: D(\mathcal A)\to H$ and satisfies
$D(\mathcal A)\hookrightarrow H^2(\Omega)\hookrightarrow C(\bar\Omega)$.

\medskip
\noindent
\textbf{Dirac actuators.}
Fix $M\in\N$ and actuator locations $x_1,\dots,x_M\in\Omega$ and define the distributions
\[
d_j := \delta_{x_j}\in D(\mathcal A)'\qquad (1\le j\le M).
\]
Define the control operator
\[
B:\R^M\to D(\mathcal A)',\qquad
Bu := \sum_{j=1}^M u_j\,d_j.
\]
This $B$ coincides with the operator in \eqref{eq:3.3} (since $D(\mathcal A)=D(A_0)$).

\begin{proposition}[Boundedness of $B$ into $D(\mathcal A)'$]
\label{prop:B-bounded-DA-dual}
There exists $C_B>0$ such that
\[
\|Bu\|_{D(\mathcal A)'} \le C_B |u|_{\R^M}\qquad \forall u\in\R^M.
\]
\end{proposition}

\begin{proof}
For $\phi\in D(\mathcal A)$ we have
$|\langle d_j,\phi\rangle| = |\phi(x_j)|\le C\|\phi\|_{H^2(\Omega)}
\le C\|\phi\|_{D(\mathcal A)}$ by Sobolev embedding and elliptic norm equivalence on
$D(\mathcal A)$. Hence
\[
|\langle Bu,\phi\rangle|
\le \sum_{j=1}^M |u_j|\,|\phi(x_j)|
\le C |u|_{\R^M}\|\phi\|_{D(\mathcal A)}.
\]
Taking the supremum over $\|\phi\|_{D(\mathcal A)}=1$ yields the claim.
\end{proof}

\medskip
\noindent
\textbf{Low-regularity solution concept.}
For $T>0$ and a right-hand side $f\in L^2(0,T;D(\mathcal A)')$ we consider
\begin{equation}
\label{eq:heat-DAdual}
\dot y(t) = A_0 y(t) + f(t)\quad\text{in }D(\mathcal A)',\qquad y(0)=y_0\in V'.
\end{equation}
A (weak/mild) solution is a function
\[
y\in W(0,T;H,D(\mathcal A)'):=\{y\in L^2(0,T;H):\dot y\in L^2(0,T;D(\mathcal A)')\}
\]
satisfying \eqref{eq:heat-DAdual} in $D(\mathcal A)'$.

\begin{proposition}[Well-posedness for Dirac-forced heat equation]
\label{prop:wp-Dirac}
\label{prop:well-posed-point}
Let $T>0$, $y_0\in V'$ and $u\in L^2(0,T;\R^M)$. Set $f=Bu\in L^2(0,T;D(\mathcal A)')$.
Then \eqref{eq:heat-DAdual} admits a unique solution $y\in W(0,T;H,D(\mathcal A)')$.
Moreover,
\begin{equation}
\label{eq:wp-estimates}
\|y\|_{L^2(0,T;H)}^2 + \|y\|_{C([0,T];V')}^2
\le C_T\Big(\|y_0\|_{V'}^2 + \|u\|_{L^2(0,T;\R^M)}^2\Big),
\end{equation}
for a constant $C_T$ depending on $T$, $\Omega$ and $\kappa_m$.
\end{proposition}

\begin{proof}
We refer to \cite{KunRodWal24-cocv}
for a detailed proof strategy in the same distribution-actuator setting.
Estimate \eqref{eq:wp-estimates} follows from the corresponding a-priori estimate. Such estimates are established under few assumptions on the domain $\Omega$, in particular it was assumed that $\Omega$ has a cubic form (or eventually convex). Such assumption is removed and replaced by a $C^{1, 1}$-regularity, see Appendix \ref{subsec:app-KRW-act-C11}. 
\end{proof}

\section{Feedback design and stability analysis}
\label{sec:feedback}
\noindent
In this section we design an output-based feedback law for the heat equation
with point actuators and show a stabilization / tracking result.
\newline
We work with the abstract system
\begin{equation}
\label{eq:feedback-abstract-system}
\dot y(t) = A_0 y(t) + B u(t),\qquad y(0)=y_0\in H,
\end{equation}
where $\Omega\subset\mathbb{R}^3$ is a bounded $C^2$ domain, In Sections~\ref{sec:feedback}--\ref{sec:coupling} we work in dimension $d=3$ and assume $\Omega$ is $C^2$ to match the thermo--plasmonic coupling; the purely abstract actuator verification in Appendix~\ref{subsec:app-KRW-act-C11} actually only requires $C^{1,1}$ regularity and applies in $d\in\{1,2,3\}$.
 $A_0$ is the
Neumann diffusion operator defined in \eqref{eq:3.2}, and $B:\mathbb{R}^M\to D(\mathcal A)'$
is the point-actuator operator defined in \eqref{eq:3.3}. We recall that $A_0$
generates an analytic contraction semigroup $(S(t))_{t\ge0}$ on $H$ and that
$B\in\mathcal{L}(\mathbb{R}^M,D(\mathcal A)')$
(Proposition~\ref{prop:B-bounded-DA-dual}).

\subsection{Closed-loop well-posedness}
\noindent
We start with the observation operator used in the feedback law.

\begin{definition}[Observation operator]
\label{def:obs-C-full}
Since $1\in\rho(A_0)$, the resolvent $\mathcal A^{-1}=(I-A_0)^{-1}$ defines a bounded operator
$H\to D(A_0)\subset H^2(\Omega)\hookrightarrow C(\overline\Omega)$. We define
\begin{equation}
\label{eq:def-C-full}
C:H\to\mathbb{R}^M,\qquad
(Cz)_j := \big(\mathcal A^{-1}z\big)(x_j),\quad j=1,\dots,M.
\end{equation}
Thus $C$ maps the state $z$ to the values of the regularized temperature
$\mathcal A^{-1}z$ at the actuator positions $x_j$.
\end{definition}
\noindent
We state
the closed-loop equation
\begin{equation}
\label{eq:closed-loop}
\dot y(t) = A_0 y(t) - \lambda B K C (y(t)-y_r),
\end{equation}
or, in terms of $z(t):=y(t)-y_r$,
\begin{equation}
\label{eq:closed-loop-z}
\dot z(t) = (A_0 - \lambda BKC) z(t),\qquad z(0)=y_0-y_r.
\end{equation}

\begin{remark}[Closed-loop well-posedness in the $V'$ framework]
\label{thm:closed-loop-wp}
For any $T>0$, $y_0\in V'$ and $u\in L^2(0,T;\R^M)$, the forced heat equation
\eqref{eq:heat-DAdual} admits a unique solution $y\in W(0,T;H,D(\mathcal A)')$ and satisfies the
estimate \eqref{eq:wp-estimates}; see Proposition~\ref{prop:wp-Dirac}.
Moreover, the feedback closed loop \eqref{eq:closed-loop-z} is well posed in $V'$.
\end{remark}
\noindent
For a given integer $M\in\mathbb{N}$ we introduce the low-mode subspace and
its projection.

\begin{definition}[Finite-dimensional subspace and projection]
\label{def:HM-PM-4}
For each $M\in\mathbb{N}$, set
\[
H_M := \mathrm{span}\{\varphi_1,\dots,\varphi_M\}\subset H,
\]
and let $P_M:H\to H_M$ be the orthogonal projection onto $H_M$. We denote
by $H_M^\perp$ the orthogonal complement of $H_M$ in $H$ and by
$I-P_M:H\to H_M^\perp$ the projection onto the tail.
\end{definition}
\noindent
We now write $C$ in the eigenbasis and separate the contributions of the low
and high modes.

\begin{lemma}[Matrix representation of $C$ and low/high decomposition]
\label{lem:C-matrix}
Let $z=\sum_{k\ge1}\alpha_k\varphi_k\in H$. Then
\[
\mathcal A^{-1}z = \sum_{k\ge1}\frac{\alpha_k}{1+\lambda_k}\,\varphi_k,
\]
and, for each $j=1,\dots,M$,
\[
(Cz)_j = \sum_{k\ge1} d_{jk}\,\alpha_k,\qquad
d_{jk} := \frac{1}{1+\lambda_k}\,\varphi_k(x_j).
\]
Define the (infinite) matrix $D=(d_{jk})_{1\le j\le M,\ k\ge1}$. Then
\[
Cz = D\alpha,\qquad \alpha:=(\alpha_1,\alpha_2,\dots)^\top\in\ell^2.
\]
For the fixed $M$ (number of actuators and low modes), decompose
$D=(D_{MM}\;|\;D_{M>})$ with
\[
D_{MM} := (d_{jk})_{j,k=1}^M\in\mathbb{R}^{M\times M},\qquad
D_{M>} := (d_{jk})_{1\le j\le M,\ k>M}.
\]
Writing $\alpha=(\alpha_M,\alpha_>)$ with
\[
\alpha_M := (\alpha_1,\dots,\alpha_M)^\top,\qquad
\alpha_> := (\alpha_{M+1},\alpha_{M+2},\dots)^\top,
\]
we obtain
\begin{equation}
\label{eq:Cz-split}
Cz = D_{MM}\alpha_M + D_{M>}\alpha_>.
\end{equation}
\end{lemma}

\begin{remark}
\label{rem:DMM-invertible}
The matrix
\[
\mathcal{B} := (\varphi_k(x_j))_{j,k=1}^M
\]
appearing in Assumption~\ref{ass:rank-4} below is related to $D_{MM}$ by
\[
D_{MM} = \mathcal{B}\,\mathrm{diag}\big((1+\lambda_k)^{-1}\big)_{k=1}^M.
\]
Hence if $\mathcal{B}$ is invertible then $D_{MM}$ is also invertible.
\end{remark}
\noindent
We use the compactness of the resolvent and the embedding $H^2(\Omega)\hookrightarrow C(\overline\Omega)$ to show that the contribution of the high modes to the measurement $Cz$ can be made arbitrarily small by choosing $N$ large enough.

\begin{lemma}[Compactness of the measurement operator and small tail]
\label{lem:tail-small-compact}
Let $\Omega\subset\mathbb{R}^3$ be a bounded $C^2$ domain, $A_0$ the Neumann
Laplacian as in \eqref{eq:3.2}, and $\{\varphi_k\}_{k\ge1}$ an
$L^2(\Omega)$–orthonormal basis of eigenfunctions of $-A_0$ with eigenvalues
$0=\lambda_1\le\lambda_2\le\cdots$. Let $x_1,\dots,x_M\in\Omega$ be fixed
distinct points, and define $C$ as in \eqref{eq:def-C-full}.
\newline
For $\alpha=(\alpha_k)_{k\ge1}\in\ell^2$, set $z=\sum_{k\ge1}\alpha_k\varphi_k$
and write $Cz=D\alpha$ with $D=(d_{jk})_{j,k}$ as in
Lemma~\ref{lem:C-matrix}. For each $N\in\mathbb{N}$, decompose
$D=(D_{MN}\mid D_{M> N})$ where
\[
D_{MN} := (d_{jk})_{1\le j\le M,\ 1\le k\le N}\in\mathbb{R}^{M\times N},\quad
D_{M> N} := (d_{jk})_{1\le j\le M,\ k>N}.
\]
\noindent
Then the operator $T:\ell^2\to\mathbb{R}^M$ defined by $T\alpha:=D\alpha$ is
compact. In particular,
\begin{equation}
\label{eq:Dtail-norm->0}
\|D_{M> N}\|_{\mathcal{L}(\ell^2,\mathbb{R}^M)} \xrightarrow[N\to\infty]{} 0.
\end{equation}
Equivalently, for every $\varepsilon>0$ there exists $N_\varepsilon\in\mathbb{N}$
such that for all $N\ge N_\varepsilon$ and all
$\alpha=(\alpha_k)_{k\ge1}\in\ell^2$,
\[
\big\|D_{M> N}\alpha_{>N}\big\|_{\mathbb{R}^M}
= \big\|C\Big( \sum_{k>N}\alpha_k\varphi_k \Big)\big\|_{\mathbb{R}^M}
\le \varepsilon\,\|\alpha\|_{\ell^2},
\]
where $\alpha_{>N} := (\alpha_{N+1},\alpha_{N+2},\dots)$.
In particular, by choosing the truncation index $N$ sufficiently large, the
measurement of the high-frequency tail can be made arbitrarily small.
\end{lemma}

\begin{proof}
Define
\[
J:\ell^2 \to H,\qquad J\alpha := \sum_{k\ge1}\alpha_k\varphi_k.
\]
Since $\{\varphi_k\}$ is an orthonormal basis of $H=L^2(\Omega)$, $J$ is an
isometric isomorphism from $\ell^2$ onto $H$.
\noindent
Let $R:=\mathcal A^{-1}:H\to D(A_0)\subset H^2(\Omega)$. Since
$H^2(\Omega)$ embeds continuously and compactly into $C(\overline\Omega)$,
then the composite operator
\[
R_J := R\circ J:\ell^2\to C(\overline\Omega)
\]
is compact. The evaluation map
\[
E: C(\overline\Omega)\to\mathbb{R}^M,\qquad
(Ew)_j := w(x_j),\quad j=1,\dots,M,
\]
is bounded. Hence the measurement operator
\[
T := E\circ R_J:\ell^2\to\mathbb{R}^M
\]
is compact as a composition of a compact and a bounded operator.
\newline
By construction, for $\alpha\in\ell^2$ and
$z=J\alpha=\sum_{k\ge1}\alpha_k\varphi_k$,
\[
T\alpha = E(RJ\alpha) = C z = D\alpha,
\]
so $T$ is represented by the matrix $D$ in the basis $\{\varphi_k\}$.
\newline
Let $P_N:\ell^2\to\ell^2$ be the orthogonal projection onto the first $N$
coordinates:
\[
P_N\alpha := (\alpha_1,\dots,\alpha_N,0,0,\dots).
\]
Set $T_N := T P_N$, which has finite rank. Since $T$ is compact, we have
\[
\|T - T_N\|_{\mathcal{L}(\ell^2,\mathbb{R}^M)}
\xrightarrow[N\to\infty]{} 0.
\]
On the other hand,
\[
(T-T_N)\alpha = T(I-P_N)\alpha,
\]
and $(I-P_N)\alpha$ has coordinates
$(0,\dots,0,\alpha_{N+1},\alpha_{N+2},\dots)$, so
\[
(T-T_N)\alpha = D_{M> N}\,\alpha_{>N}.
\]
Hence
\[
\|T-T_N\|_{\mathcal{L}(\ell^2,\mathbb{R}^M)}
= \|D_{M> N}\|_{\mathcal{L}(\ell^2,\mathbb{R}^M)},
\]
which proves \eqref{eq:Dtail-norm->0}. The last statement is just the
operator norm inequality written explicitly.
\end{proof}
\noindent
In what follows, the number $M$ denotes the number of actuators/sensors (i.e.\ the dimension
of the control and observation space $\mathbb{R}^M$). For the spectral splitting we introduce
an independent truncation index $N\in\mathbb{N}$ (typically $N\le M$). We set
\[
H_N := \mathrm{span}\{\varphi_1,\dots,\varphi_N\}
\]
and $P_N:H\to H_N$ as in Definition~\ref{def:HM-PM-4}, and we treat $(I-P_N)H$ as the tail.
Lemma~\ref{lem:tail-small-compact} shows that
\[
\|C(I-P_N)\|_{\mathcal{L}(H,\mathbb{R}^M)} \longrightarrow 0
\quad\text{as } N\to\infty;
\]
hence, for any prescribed $\delta>0$, we can choose $N$ so that
\[
\|C(I-P_N)\|_{\mathcal{L}(H,\mathbb{R}^M)} \le \delta.
\]

\subsection{Assumptions and feedback law}
\label{subsec:assumptions-4}
\noindent
We now state the structural assumptions needed for the stabilization result.

\begin{assumption}[Finite-dimensional subspace and rank condition]\label{ass:rank-4}
Fix $M\in\mathbb N$ and distinct actuator/sensor locations $x_1,\dots,x_M\in\Omega$.
Let $N\in\mathbb N$ be a truncation index (typically $N\le M$) and let $H_N$ and $P_N$
be as in Definition~\ref{def:HM-PM-4}.

\begin{enumerate}
\item[(i)] \textbf{(Rank condition on the first $N$ modes).}
The $N\times M$ matrix
\[
\Phi_{N,M}:=\big(\varphi_k(x_j)\big)_{\substack{1\le k\le N\\ 1\le j\le M}}
\]
has full row rank $N$.
Equivalently, the $M\times N$ matrix
\[
D_{M,N}:=\Big(\frac{\varphi_k(x_j)}{1+\lambda_k}\Big)_{\substack{1\le j\le M\\ 1\le k\le N}}
\]
has full column rank $N$.

\item[(ii)] \textbf{(Reference in $H_N$).}
The reference profile satisfies $y_r\in H_N$, i.e.
\[
y_r=\sum_{k=1}^N \alpha_k^r\,\varphi_k
\quad\text{for some }\alpha_k^r\in\mathbb R.
\]
\end{enumerate}
\end{assumption}
\noindent
Assumption~\ref{ass:rank-4}(i) is a controllability-type condition on the
configuration of actuator positions $x_j$ with respect to the first $N$
eigenfunctions $\varphi_k$. As discussed in the appendix (see the genericity and 1D/3D
examples), this condition is generic and can be enforced explicitly in simple
geometries.
\newline
The feedback law uses the full observation operator $C$ as defined in
\eqref{eq:def-C-full}.

\begin{definition}[Feedback law]
\label{def:feedback-law-4}
Let $K\in\mathcal{L}(\mathbb{R}^M,\mathbb{R}^M)$ be a feedback gain. We
consider the static output feedback
\begin{equation}
\label{eq:feedback-fullC}
u(t) = -K\,C\big(y(t)-y_r\big).
\end{equation}
The closed-loop system for $y$ is then
\begin{equation}
\label{eq:closed-loop-y-4}
\dot y(t) = A_0 y(t) - B K C\big(y(t)-y_r\big).
\end{equation}
For the error variable $z(t):=y(t)-y_r$ we have
\begin{equation}
\label{eq:closed-loop-z-4}
\dot z(t) = A_0 z(t) - B K C z(t),\qquad z(0)=z_0:=y_0-y_r.
\end{equation}
\end{definition}
\noindent
Well-posedness of \eqref{eq:closed-loop-z-4} follows as in
Proposition~\ref{prop:wp-Dirac}, since $BKC\in\mathcal{L}(H,H_{-1})$.

%\subsection{Detailed proof of the stabilization theorem}\label{subsec:4.3}

\begin{remark}[Finite-dimensional constructions versus stabilization]
\label{rem:finite-vs-krw}
A matrix-gain construction based on truncation to the first $N$ eigenmodes can be used as a
\emph{heuristic} guideline for choosing actuator locations and tuning feedback parameters.
However, for Dirac actuators the rigorous closed-loop well-posedness and stabilization/tracking
results are most naturally formulated in the $V'$ framework. In this paper, the main stability
statement is Theorem~\ref{thm:KRW-Vprime}, which yields exponential decay in $V'$ for the
feedback \eqref{eq:KRW-feedback} under Assumption~\ref{assump:KRW-act}, and implies $H$--tracking
after any positive time $t_0>0$ (see Corollary~\ref{cor:L2-after-t0} and Section~\ref{sec:robust-Vprime}).
\end{remark}

%==================== sec4_stabilization_KRW.tex ====================
\section{Dirac-actuator stabilization and tracking in $V'$}
\label{sec:stab-KRW}
\noindent
We rewrite the tracking problem in terms of the error variable
\[
z(t):=y(t)-y_r.
\]
Throughout this section we allow non-constant reference profiles by assuming that $y_r$ is a finite
linear combination of Neumann eigenfunctions. More precisely, for some $N\in\N$,
\begin{equation}
\label{eq:ref-finite-modes}
y_r\in X_N:=\mathrm{span}\{\varphi_1,\dots,\varphi_N\},\qquad
y_r=\sum_{k=1}^N \alpha_k^r\,\varphi_k.
\end{equation}
Since such a $y_r$ is not, in general, an equilibrium of the free heat dynamics, we introduce a constant
feedforward input $u_r\in\R^M$ chosen so that the forcing $A_0y_r+Bu_r$ is eliminated at the level of
the first $N$ modes:
\begin{equation}
\label{eq:feedforward-moment}
\langle A_0 y_r + Bu_r,\varphi_k\rangle_{V',V}=0,\qquad k=1,\dots,N.
\end{equation}
Equivalently,
\begin{equation}
\label{eq:feedforward-linear-system}
\sum_{j=1}^M (u_r)_j\,\varphi_k(x_j) = -\langle A_0 y_r,\varphi_k\rangle_H,\qquad k=1,\dots,N.
\end{equation}
We assume that the $N\times M$ matrix $(\varphi_k(x_j))_{k=1,\dots,N;\ j=1,\dots,M}$ has full row rank,
so that \eqref{eq:feedforward-linear-system} is solvable; when $M>N$ we take $u_r$ to be the minimum-norm
solution of this system.

\subsection{Feedback law}
\noindent
Let $\mathcal A=I-A_0$ and $d_j=\delta_{x_j}\in D(\mathcal A)'$ as in Section~\ref{subsec:wp-Dirac-KRW}.
Define the (elliptically smoothed) point measurements
\[
\big(Cz\big)_j := \langle d_j,\mathcal A^{-1}z\rangle_{D(\mathcal A)',D(\mathcal A)}
= \big(\mathcal A^{-1}z\big)(x_j),\qquad j=1,\dots,M.
\]
We consider the feedback
\begin{equation}
\label{eq:KRW-feedback}
u_j(t) = (u_r)_j -\lambda\, (Cz(t))_j,\qquad j=1,\dots,M,
\end{equation}
with a scalar gain $\lambda>0$. Equivalently, the closed-loop error dynamics reads
\begin{equation}
\label{eq:KRW-closed-loop}
\dot z(t)=A_0 z(t) -\lambda\sum_{j=1}^M \langle d_j,\mathcal A^{-1}z(t)\rangle\, d_j + g
\quad\text{in }D(\mathcal A)',\qquad z(0)=z_0\in V',
\end{equation}
where the (time-independent) forcing term is
\begin{equation}
\label{eq:feedforward-g}
g:=A_0 y_r + B u_r\in D(\mathcal A)'.
\end{equation}
\noindent
Observe that, if we are only interested in equilibrium states, which means $A_0y_r=0$ (and $\partial_n y_r=0$ on $\partial \Omega$), then we take $u_r=0$ so that $g=0$.
\subsection{Actuator family condition}
\noindent
The stabilization mechanism requires a geometric/approximation condition on the actuators.
We adopt the abstract actuator hypothesis of \cite{KunRodWal24-cocv}
(Assumption~2.4 therein), specialized to Dirac actuators.

\begin{assumption}[Actuator hypothesis \cite{KunRodWal24-cocv}]
\label{assump:KRW-act}
There exists an increasing sequence $M\mapsto M_\sigma$ and, for each $M$, a family of
actuators $\{d_{M,j}\}_{j=1}^{M_\sigma}\subset D(\mathcal A)'$ and auxiliary functions
$\{\Psi_{M,j}\}_{j=1}^{M_\sigma}\subset D(\mathcal A)$ such that:
\begin{enumerate}
\item $\dim\mathrm{span}\{d_{M,j}\}=M_\sigma$ and $\dim\mathrm{span}\{\Psi_{M,j}\}=M_\sigma$;
\item biorthogonality: $\langle d_{M,j},\Psi_{M,i}\rangle=\delta_{ij}$;
\item the ``nullspace coercivity'' constant
\[
\xi_M^+ := \inf\Big\{\frac{\|w\|_{D(\mathcal A)}^2}{\|w\|_{V}^2}\,:\,
w\in D(\mathcal A)\setminus\{0\},\ \langle d_{M,j},w\rangle=0\ \forall j\Big\}
\]
satisfies $\xi_M^+\to+\infty$ as $M\to\infty$.
\end{enumerate}
\end{assumption}

\begin{remark}
In box (or more general polyhedral domains) domains $\Omega$ and for appropriately distributed actuator points, we can follow~\cite{KunRodWal24-cocv}
to verify Assumption~\ref{assump:KRW-act} (Section~5 therein) using self-similar domain decompositions.
In Appendix~\ref{subsec:app-KRW-act-C11} we show that, if $\Omega$ is a bounded $C^{1,1}$ domain (hence in particular
for the $C^2$ domains assumed in the main text), then Assumption~\ref{assump:KRW-act} can be verified for Dirac actuators
placed at the vertices of any shape-regular quasi-uniform triangulation of $\Omega$. In particular,
the nullspace coercivity constant satisfies $\xi_M^+\gtrsim h_M^{-2}\to\infty$ as the mesh size $h_M\downarrow 0$ (equivalently
$\xi_M^+\gtrsim M_\sigma^{2/d}$).
\end{remark}

\subsection{Exponential tracking in $V'$ and an $L^2$ consequence for $t\ge t_0$}

\begin{theorem}[Exponential stabilization in $V'$ and tracking]
\label{thm:KRW-Vprime}
Assume Assumption~\ref{assump:KRW-act}. For every decay rate $\mu>0$ there exist $M$ large enough
and $\lambda>0$ large enough such that the \emph{homogeneous} closed loop (i.e.\ \eqref{eq:KRW-closed-loop}
with $g=0$) generates an exponentially stable semigroup $\mathcal{S}_{\mathrm{cl}}(t)$ on $V'$ with
$\|\mathcal{S}_{\mathrm{cl}}(t)\|_{\mathcal{L}(V')} \le e^{-\mu t}$. Consequently, for the affine closed loop
\eqref{eq:KRW-closed-loop} there exists an equilibrium $z_\infty\in V'$ given by
\begin{equation}
\label{eq:z-infty}
z_\infty := \int_0^\infty \mathcal{S}_{\mathrm{cl}}(s)\,g\,\dd s,
\end{equation}
and the corresponding solution $z$ satisfies
\begin{equation}
\label{eq:Vprime-decay}
\|z(t)-z_\infty\|_{V'} \le e^{-\mu (t-s)} \|z(s)-z_\infty\|_{V'}\qquad \forall t\ge s\ge 0.
\end{equation}
Moreover, there exists $C_\mu>0$ such that
\begin{equation}
\label{eq:L2-time-bound}
\|z-z_\infty\|_{L^2((s,\infty);H)} \le C_\mu \|z(s)-z_\infty\|_{V'}\qquad \forall s\ge 0.
\end{equation}
\end{theorem}

\begin{proof}
For the homogeneous closed loop ($g=0$), this is a direct specialization of
\cite[Theorem~3.1]{KunRodWal24-cocv} to the autonomous case
($A_{rc}\equiv 0$) with $\mathcal A=I-A_0$ and Dirac actuators, which yields
$\|\mathcal{S}_{\mathrm{cl}}(t)\|_{\mathcal{L}(V')}\le e^{-\mu t}$.
\noindent
For the affine system \eqref{eq:KRW-closed-loop}, the variation-of-constants formula gives
\[
z(t)=\mathcal{S}_{\mathrm{cl}}(t)z_0+\int_0^t \mathcal{S}_{\mathrm{cl}}(t-s)\,g\,\dd s.
\]
Since $\mathcal{S}_{\mathrm{cl}}$ is exponentially stable on $V'$, the improper integral in
\eqref{eq:z-infty} converges in $V'$ and defines an equilibrium $z_\infty$.
Moreover,
\[
z(t)-z_\infty=\mathcal{S}_{\mathrm{cl}}(t)\,(z_0-z_\infty),
\]
which implies \eqref{eq:Vprime-decay}. The $L^2$ bound \eqref{eq:L2-time-bound} follows from the
corresponding estimate applied to the homogeneous
trajectory $z-z_\infty$.
\end{proof}

\begin{corollary}[$L^2$ tracking for $t\ge t_0>0$]
\label{cor:L2-after-t0}
Let the assumptions of Theorem~\ref{thm:KRW-Vprime} hold and fix $t_0>0$.
Then there exists $C(t_0)>0$ such that
\begin{equation}
\label{eq:L2-after-t0}
\sup_{t\ge t_0}\|z(t)-z_\infty\|_{H}
\le C(t_0)\,\|z-z_\infty\|_{L^2((t_0/2,\infty);H)}.
\end{equation}
In particular, combining \eqref{eq:L2-time-bound} with \eqref{eq:L2-after-t0} (with $s=t_0/2$) yields
\[
\sup_{t\ge t_0}\|z(t)-z_\infty\|_{L^2(\Omega)}
\le C(t_0)\,C_\mu\, e^{-\mu t_0/2}\,\|z_0-z_\infty\|_{V'}.
\]
\end{corollary}

\begin{proof}
The bound \eqref{eq:L2-after-t0} is a standard parabolic smoothing estimate; see, e.g., \cite[Ch.~4]{Pazy1983} or
\cite[Ch.~2]{Lunardi-Book}. One can also prove it by expanding in the eigenbasis of $-A_0$ and using that high modes
are exponentially damped on any interval of length $t_0/2$.
\end{proof}

\begin{remark}[What is (and is not) claimed]
\label{rem:KRW-what-claimed}
Theorem~\ref{thm:KRW-Vprime} is the main rigorous tracking statement (in $V'$).
Corollary~\ref{cor:L2-after-t0} provides an $L^2$ consequence on $[t_0,\infty)$, with a constant
depending on $t_0$. It does not claim a uniform $t_0\downarrow 0$ bound.
\end{remark}

%--------------------------------------------------------
\subsection{Fixed-point pre-compensation of the steady bias}
\label{subsec:fixed-point-bias}
%--------------------------------------------------------
\noindent
The discussion above shows that for a non-equilibrium reference profile $y_r$ the closed-loop error
$z(t):=y(t)-y_r$ converges to a (possibly nonzero) equilibrium
$z_\infty$ given by \eqref{eq:z-infty}. This subsection provides a systematic way to compensate for this
bias by adjusting the \emph{commanded} reference within the finite-dimensional space $X_N$ introduced in
\eqref{eq:ref-finite-modes}. The construction is based on a fixed-point equation.

\medskip
\noindent\textbf{The bias map.}
Fix $N\in\mathbb{N}$ and recall $X_N:=\mathrm{span}\{\varphi_1,\dots,\varphi_N\}$.
For any commanded reference $y^\star\in X_N$ we define the feedforward $u_r(y^\star)\in\mathbb{R}^M$ as a
(possibly nonunique) solution of the moment system \eqref{eq:feedforward-linear-system}--\eqref{eq:feedforward-moment}
with $y_r$ replaced by $y^\star$. In what follows we fix a \emph{linear} selection
\begin{equation}
\label{eq:UN-def}
U_N:X_N\to\mathbb{R}^M,\qquad u_r(y^\star):=U_N y^\star,
\end{equation}
for instance the minimum--Euclidean-norm choice, which is linear whenever $\Phi_{N,M}$ has full row rank.
With this choice we associate to $y^\star$ the residual (time-independent) forcing
\begin{equation}
\label{eq:g-of-ystar}
g(y^\star):=A_0 y^\star + B\,U_N y^\star \in D(\mathcal{A})'.
\end{equation}
By construction we have $P_N g(y^\star)=0$ (cancellation of the first $N$ moments), and since
$A_0y^\star\in X_N$ this equivalently reads
\begin{equation}
\label{eq:g-is-tail}
g(y^\star)=(I-P_N)\,B\,U_N y^\star .
\end{equation}
\noindent
Let $\mathcal{S}_{\mathrm{cl}}(t)$ denote the exponentially stable closed-loop semigroup associated with the
homogeneous loop on $V'$ (Theorem~\ref{thm:KRW-Vprime}). As used in \eqref{eq:z-infty}, the affine loop with
forcing $g\in D(\mathcal A)'$ admits an equilibrium $z_\infty\in V'$ given by the convergent integral
$\int_0^\infty \mathcal{S}_{\mathrm{cl}}(s)\,g\,\dd s$; see \cite{KunRodWal24-cocv}.
Accordingly, we define the (bounded) operator
\begin{equation}
\label{eq:Rcl-def}
\mathcal{R}_{\mathrm{cl}}:D(\mathcal A)'\to V',\qquad
\mathcal{R}_{\mathrm{cl}}h := \int_0^\infty \mathcal{S}_{\mathrm{cl}}(s)\,h\,\dd s,
\end{equation}
and denote by $C_{\mathrm{cl}}>0$ any constant such that
\begin{equation}
\label{eq:Rcl-bound}
\|\mathcal{R}_{\mathrm{cl}}h\|_{V'} \le C_{\mathrm{cl}}\,\|h\|_{D(\mathcal A)'}\qquad\forall h\in D(\mathcal A)'.
\end{equation}

\begin{lemma}[Linearity and a priori bound for the steady bias]
\label{lem:bias-map}
The map
\begin{equation}
\label{eq:Zmap-def}
\mathcal{Z}_N:X_N\to V',\qquad
\mathcal{Z}_N(y^\star):=\mathcal{R}_{\mathrm{cl}}\,g(y^\star),
\end{equation}
is linear and bounded. Moreover, for all $y^\star\in X_N$,
\begin{equation}
\label{eq:Zmap-bound}
\|\mathcal{Z}_N(y^\star)\|_{V'} \le C_{\mathrm{cl}}\,\|g(y^\star)\|_{D(\mathcal A)'}.
\end{equation}
\end{lemma}

\begin{proof}
Linearity follows from the linearity of $U_N$ and $B$ in \eqref{eq:g-of-ystar} and the linearity of
$\mathcal{R}_{\mathrm{cl}}$ in \eqref{eq:Rcl-def}. The bound \eqref{eq:Zmap-bound} is \eqref{eq:Rcl-bound}.
\end{proof}

\begin{remark}[Achieved steady field for a commanded reference]
\label{rem:achieved-steady-field}
For a commanded reference $y^\star\in X_N$ we consider the shifted error $z:=y-y^\star$ and apply the
feedback \eqref{eq:KRW-feedback} with the corresponding feedforward $u_r(y^\star)=U_Ny^\star$.
Then the closed-loop error is affine with forcing $g(y^\star)$ and converges exponentially in $V'$ to the equilibrium
\[
z_\infty(y^\star)=\mathcal{Z}_N(y^\star),
\]
so that the temperature field converges to
\begin{equation}
\label{eq:yinfty-of-ystar}
y_\infty(y^\star):=\lim_{t\to\infty} y(t) = y^\star + \mathcal{Z}_N(y^\star)\qquad \text{in }V'.
\end{equation}
\end{remark}

\medskip
\noindent\textbf{Fixed-point pre-compensation in $X_N$.}
Given a desired target $y_r\in X_N$, we seek a commanded reference $y^\star\in X_N$ such that the achieved
steady field \eqref{eq:yinfty-of-ystar} matches $y_r$ \emph{on the low modes}. Since $X_N\subset V$ we can use the
$V'$--$V$ duality pairing to define the spectral projector $P_N$ on $V'$ and $D(\mathcal A)'$ by
\[
P_N \eta := \sum_{k=1}^N \langle \eta,\varphi_k\rangle_{D(\mathcal A)',D(\mathcal A)}\,\varphi_k,
\qquad \eta\in D(\mathcal A)'.
\]
Define the projected bias operator
\begin{equation}
\label{eq:TN-def}
T_N:X_N\to X_N,\qquad T_N y := P_N \mathcal{Z}_N(y).
\end{equation}
The compensated target is then determined by the linear fixed-point equation
\begin{equation}
\label{eq:fixed-point}
y^\star + T_N y^\star = y_r\qquad\text{in }X_N,
\end{equation}
i.e.\ $(I+T_N)y^\star=y_r$.

\begin{proposition}[Existence/uniqueness and iterative construction]
\label{prop:fixed-point}
Let $y_r\in X_N$.
\begin{enumerate}
\item If $I+T_N:X_N\to X_N$ is invertible, then \eqref{eq:fixed-point} admits a unique solution
$y^\star\in X_N$. For this choice of commanded reference one has
\begin{equation}
\label{eq:low-mode-matching}
P_N y_\infty(y^\star)=y_r.
\end{equation}
\item If, in addition, $\|T_N\|_{\mathcal{L}(X_N)}<1$ in some norm on $X_N$, then the Picard iteration
\begin{equation}
\label{eq:Picard}
y^{(0)}:=y_r,\qquad y^{(\ell+1)}:=y_r- T_N y^{(\ell)},\qquad \ell\ge 0,
\end{equation}
converges geometrically in $X_N$ to the unique solution $y^\star$ of \eqref{eq:fixed-point}.
\end{enumerate}
\end{proposition}

\begin{proof}
(1) Since $X_N$ is finite-dimensional, invertibility of $I+T_N$ is equivalent to existence and uniqueness of the
solution $y^\star=(I+T_N)^{-1}y_r$. The identity \eqref{eq:low-mode-matching} follows from
\eqref{eq:yinfty-of-ystar} and \eqref{eq:fixed-point}:
\[
P_N y_\infty(y^\star)=P_N\bigl(y^\star+\mathcal{Z}_N(y^\star)\bigr)=y^\star+T_Ny^\star=y_r.
\]
(2) The iteration \eqref{eq:Picard} is the standard Banach fixed-point iteration for the contraction
$F(y):=y_r-T_Ny$ on $X_N$.
\end{proof}

\begin{corollary}[Exact low-mode tracking and a quantitative tail estimate]
\label{cor:fixed-point-tail}
Assume $y_r\in X_N$ and let $y^\star\in X_N$ be any solution of the fixed-point equation \eqref{eq:fixed-point}.
Define the corresponding steady reached profile by \eqref{eq:yinfty-of-ystar}.
Then
\begin{equation}
\label{eq:exact-low-mode-tracking}
P_N y_\infty(y^\star)=y_r,
\end{equation}
and the full-field mismatch is a pure tail:
\begin{equation}
\label{eq:tail-only-mismatch}
y_\infty(y^\star)-y_r = (I-P_N)\,y_\infty(y^\star)=(I-P_N)\,\mathcal{Z}_N(y^\star).
\end{equation}
In particular,
\begin{equation}
\label{eq:tail-estimate-Vprime}
\|y_\infty(y^\star)-y_r\|_{V'}
\le \|I-P_N\|\,C_{\mathrm{cl}}\,
\|(I-P_N)B\|_{\mathcal{L}(\mathbb{R}^M,D(\mathcal A)')}\,\|U_N\|\,\|y^\star\|_{X_N}.
\end{equation}
If $I+T_N$ is invertible, so that $y^\star=(I+T_N)^{-1}y_r$, then \eqref{eq:tail-estimate-Vprime} yields
\begin{equation}
\label{eq:tail-estimate-with-inverse}
\|y_\infty(y^\star)-y_r\|_{V'}
\le \|I-P_N\|\,C_{\mathrm{cl}}\,
\|(I-P_N)B\|\,\|U_N\|\,\|(I+T_N)^{-1}\|\,\|y_r\|_{X_N}.
\end{equation}
Moreover, if $\|T_N\|_{\mathcal{L}(X_N)}<1$, then $\|(I+T_N)^{-1}\|\le (1-\|T_N\|)^{-1}$.
\end{corollary}

\begin{proof}
The low-mode identity \eqref{eq:exact-low-mode-tracking} is \eqref{eq:low-mode-matching}.
Since $y^\star\in X_N$, we have $(I-P_N)y^\star=0$, and therefore by \eqref{eq:yinfty-of-ystar}
\[
(I-P_N)\bigl(y_\infty(y^\star)-y_r\bigr)
=(I-P_N)\bigl(y^\star+\mathcal{Z}_N(y^\star)-y_r\bigr)
=(I-P_N)\mathcal{Z}_N(y^\star),
\]
which is \eqref{eq:tail-only-mismatch}. Next, combining \eqref{eq:Zmap-bound} with \eqref{eq:g-is-tail} gives
\[
\|y_\infty(y^\star)-y_r\|_{V'}
=\|(I-P_N)\mathcal{Z}_N(y^\star)\|_{V'}
\le \|I-P_N\|\,\|\mathcal{Z}_N(y^\star)\|_{V'}
\le \|I-P_N\|\,C_{\mathrm{cl}}\,\|(I-P_N)B\|\,\|U_N\|\,\|y^\star\|_{X_N},
\]
which is \eqref{eq:tail-estimate-Vprime}. The bound \eqref{eq:tail-estimate-with-inverse} follows from
$y^\star=(I+T_N)^{-1}y_r$.
\end{proof}


\begin{remark}[A sufficient small-gain condition]
\label{rem:small-gain-TN}
From \eqref{eq:TN-def}, Lemma~\ref{lem:bias-map}, and \eqref{eq:g-is-tail} one obtains the bound
\[
\|T_N\|_{\mathcal{L}(X_N)}
\;\le\; \|P_N\|\,C_{\mathrm{cl}}\,\|(I-P_N)B\|_{\mathcal{L}(\mathbb{R}^M,D(\mathcal A)')}\,\|U_N\|.
\]
Hence $\|T_N\|<1$ follows if the tail $\|(I-P_N)B\|$ is sufficiently small and the moment
solver $U_N$ is not too ill-conditioned (in particular, if $\Phi_{N,M}$ has a uniformly positive smallest singular value).
\end{remark}

\begin{remark}[How to compute $T_N$ in practice]
\label{rem:compute-TN}
Represent $y\in X_N$ by its coefficient vector $a\in\mathbb{R}^N$ in the basis $\{\varphi_k\}_{k=1}^N$.
To assemble the matrix of $T_N$, for each $k=1,\dots,N$ set $y^\star=\varphi_k$, compute $u_r=U_N\varphi_k$,
form $g(\varphi_k)$ via \eqref{eq:g-of-ystar}, and compute the equilibrium $z_\infty(\varphi_k)=\mathcal{Z}_N(\varphi_k)$
(by \eqref{eq:Rcl-def}) through the stationary closed-loop problem $A_{\mathrm{cl}}z+g(\varphi_k)=0$. Then the $(\ell,k)$ entry of the
$N\times N$ matrix of $T_N$ is
\[
(T_N)_{\ell k} = \langle z_\infty(\varphi_k),\varphi_\ell\rangle_{V',V},\qquad 1\le \ell,k\le N,
\]
and \eqref{eq:fixed-point} becomes the linear system $(I+T_N)a^\star=a^r$.
\end{remark}



\section{Coupling with the thermo-plasmonic model and tracking}
\label{sec:coupling}
\noindent
We now explain how the abstract control model of Sections~\ref{sec:point-actuators}
and \ref{sec:feedback} is realized by the discrete thermo--plasmonic model
of Section~\ref{sec:model}, and we derive a tracking result for the
nanoparticle-based actuator system.

\subsection{From the whole-space model to a bounded control domain}
\noindent
The effective thermo--plasmonic model is posed in $\R^3$
(Theorem~\ref{thm:discrete-effective}), while the control analysis is
performed on a bounded domain $\Om$. We briefly justify this reduction.

\begin{assumption}[Geometric setting for the control domain]
\label{ass:Omega-large}
Let $B\subset\R^3$ be a fixed compact set that contains all nanoparticle
centers $z_i$ and the application-relevant region (e.g.\ the tumour region
where heat tracking is required). We choose $\Om$ to be a bounded $C^2$
domain such that
\[
B \Subset \Om \Subset \R^3,
\]
and the points $x_j$ where point actuators are located satisfy $x_j\in B$.
\end{assumption}
\noindent
Under Assumption~\ref{ass:Omega-large}, the whole-space solution $w$ of the
effective model can be restricted to $\Om$, and the difference between the
whole-space solution and the bounded-domain one is small on $B$ if $\Om$
is chosen large. This reduction is quantified in Lemma \ref{lem:R3-to-Omega-restriction}.


\subsection{Abstract map from intensity controls to heat inputs} \label{sec:map-EM-Heat}
\noindent
We now summarize the abstract structure of the map from electromagnetic actuation to effective
heat inputs.
In our heat-generation model, the volumetric source is parameterized as a \textit{nonnegative} linear combination of precomputed intensity patterns,  see (\ref{eq:heat-transmission}) and (\ref{Q_i-s}),,:
\begin{equation}\label{eq:heat-source-dictionary}
q(\cdot,t)=\sum_{\ell=1}^P p_\ell(t)\,\psi_\ell(\cdot),
\qquad 
\psi_\ell(\cdot)=\bigl|E_{\mathrm{inc},\ell}(\cdot)\bigr|^2.
\end{equation}
Here each dictionary element corresponds to a choice of incident direction and polarization (and, when relevant, frequency and nanoparticle parameters near resonance). Thus, the actuation map $p\mapsto q$ is linear by construction, with physically admissible coefficients $p_\ell(t)\ge 0$. Consequently, the forcing depends only on the \emph{intensity weights},
\[
 p(t)=(p_\ell(t))_{\ell=1}^P\in\R^P,
\]
and the reduced actuation map is linear in $p$. Throughout the remainder of the paper, $p$ denotes this
(real-valued) intensity/power control parameter.
Let $U$ be a Hilbert space of time-dependent intensity controls (e.g.\ $U=H^{r+1/2}_{0,\sigma}(0,T;\R^P)$). As in
Section~\ref{sec:model}, the Volterra system \eqref{eq:Volterra-system} can be written in operator form
\[
\sigma = \mathcal{V}F + r,
\]
where $F=(F_i)$ is the forcing vector depending linearly on the intensity control $p$, $\mathcal{V}$ is a bounded Volterra operator on
$X:=(H^1(0,T))^M$, and $r$ is a small remainder. The effective heat
inputs associated with the nanoparticles are
\begin{equation}
\label{eq:Gi-sigma}
G_i(t) := \frac{\alpha_i}{c_m}\,\sigma_i(t),\qquad i=1,\dots,M,
\end{equation}
and we denote $G=(G_i)_{i=1}^M$.

\begin{theorem}[Rigorous approximation of the map $p\mapsto G$]
\label{thm:K-rigorous}
Let $U:=H^{r+1/2}_{0,\sigma}(0,T;\R^P)$ be the space of intensity control
parameters, $X:=(H^1(0,T))^M$ and $Y:=L^2(0,T;\R^M)$. Assume that:
\begin{enumerate}
\item[(i)] The map $p\mapsto F(p)$ associated
with the Maxwell problem and the forcing term in \eqref{eq:Volterra-system} extends to a bounded
linear operator $\mathcal{T}:U\to X$.

\item[(ii)] The Volterra system \eqref{eq:Volterra-system} defines a bounded
linear operator $\mathcal{V}\in\mathcal{L}(X,X)$ and the remainder
$r\in X$ satisfies
\[
\norm{r}_X \le C_T M\delta^\mu\norm{F}_X
\]
for some $\mu>0$ and $C_T>0$ independent of $\delta$.

\item[(iii)] The map $\sigma\mapsto G$ given by \eqref{eq:Gi-sigma} extends
to a bounded linear operator $\mathcal{S}\in\mathcal{L}(X,Y)$.
\end{enumerate}
Then there exist a bounded linear operator $\mathcal{K}:U\to Y$ and a
remainder operator $\mathcal{R}^\delta:U\to Y$ such that, for every
$p\in U$,
\begin{equation}
\label{eq:K-decomposition}
G = \mathcal{K}p + \mathcal{R}^\delta p,
\end{equation}
with
\begin{equation}
\label{eq:Rdelta-bound}
\norm{\mathcal{R}^\delta p}_Y \le C_T' M\delta^\mu\norm{p}_U,
\end{equation}
for some constant $C_T'>0$ independent of $\delta$ and $p$.
\end{theorem}

\begin{proof}
As before: $\sigma = \mathcal{V}F + r$ with
$\norm{r}_X\le C_T M\delta^\mu\norm{F}_X$, $F=\mathcal{T}p$, and
$G=\mathcal{S}\sigma$. Define $\mathcal{K}:=\mathcal{S}\mathcal{V}
\mathcal{T}$ and $\mathcal{R}^\delta p := \mathcal{S}r(p)$.
\end{proof}

\begin{remark}[\textbf{Constraints on the intensity control $p$ (nonnegativity)}] \label{Positivity-p}
In the physical model the illumination/intensity parameters satisfy
$p(t)=(p_\ell(t))_{\ell=1}^P\in(\mathbb{R}_+)^P$ since $p_\ell(t)=|a_\ell(t)|^2\ge 0$.
In the reduced implementation the actuator signal that enters the
Dirac-driven heat model is a finite-dimensional quantity
$u_{\mathrm{des}}(t)\in Y_0$ (obtained by projecting the abstract feedback input onto $Y_0$), and it is related to
the intensity vector through the linear map $K_0:U_0\to Y_0$.
Therefore, with the physical constraint $p_l(t)\ge 0, l=1, ... P$, the set of \emph{instantaneously implementable} actuator inputs
is the convex cone
\[
\mathcal{C}_{K_0}:=K_0(\mathbb{R}_+^P)\subset Y_0 .
\]
Consequently, exact realization of the commanded input is possible at time $t$ if and only if
$u_{\mathrm{des}}(t)\in \mathcal{C}_{K_0}$, in which case one may select (for instance) the minimal-norm nonnegative
intensity vector
\[
p(t)\in \arg\min_{p\in\mathbb{R}_+^P}\ \|p\|^2
\quad\text{subject to}\quad K_0 p = u_{\mathrm{des}}(t),
\]
for instance.
When $u_{\mathrm{des}}(t)\notin \mathcal{C}_{K_0}$, one cannot represent it exactly with $p(t)\ge 0$ using the fixed
dictionary; in that case a natural implementation is to use the \emph{cone projection}
\[
p(t)\in \arg\min_{p\in\mathbb{R}_+^P}\ \|K_0 p-u_{\mathrm{des}}(t)\|_{Y_0}^2,
\qquad u_{\mathrm{impl}}(t):=K_0 p(t)\in\mathcal{C}_{K_0},
\]
so that the closed-loop system tracks the best attainable input within $\mathcal{C}_{K_0}$.
In the remainder, $p$ denotes a physically admissible (nonnegative) intensity vector produced by the above procedures. Such arguments will be exploited in a forthcoming work.
\end{remark}







\begin{proposition}[Verification of (i)--(iii) from the discrete model]
\label{prop:verification-K}
Under the assumptions and notation of Theorem~\ref{thm:discrete-effective},
the hypotheses (i)--(iii) of Theorem~\ref{thm:K-rigorous} are satisfied,
with $\mathcal{T},\mathcal{V},\mathcal{S}$ constructed from the discrete
effective thermo--plasmonic model of \cite{CaoMukherjeeSini2025}. In
particular, the approximation \eqref{eq:K-decomposition} holds for all
intensity controls $p\in U$.
\end{proposition}

\begin{proof}
The arguments are as follows: (i) uses the linear dependence
of $F_i$ on the EM fields, (ii) follows from the invertibility of
$I+\mathcal{W}$ in the Volterra system and the error estimates in
\cite{CaoMukherjeeSini2025}, and (iii) is immediate from
\eqref{eq:Gi-sigma} and the embedding $H^1(0,T)\hookrightarrow L^2(0,T)$.
\end{proof}

\subsection{Approximate right-inverse for \texorpdfstring{$\mathcal{K}$}{K}}\label{sec:K-inverse}
\noindent
We now state conditions under which $\mathcal{K}$ admits an approximate
right-inverse on a finite-dimensional subspace of inputs, as used in the
tracking theorem.

\begin{assumption}[Non-degenerate illumination patterns]
\label{ass:illumination}
Let $D_i=z_i+\delta B_i$, $i=1,\dots,M$, be the nanoparticle configuration
satisfying the assumptions of Theorem~\ref{thm:discrete-effective}, and fix
a finite family of incident fields $E^{(\ell)}_{\mathrm{in}}$,
$\ell=1,\dots,P$, with corresponding total fields $E^{(\ell)}$ solving
\eqref{eq:maxwell}. Let $c_i^{(\ell)}(\delta)$ denote the leading-order
coefficient relating $|E^{(\ell)}|^2$ to the forcing term $F_i$ in
\eqref{eq:Volterra-system}. Assume that
\[
c_i^{(\ell)}(\delta) = c_i^{(\ell)}(0) + \mathcal{O}(\delta^\mu),
\]
and that the $M\times P$ matrix $L_0=(c_i^{(\ell)}(0))$ has rank $M$, with
a uniformly bounded right-inverse on its range. In particular, $P\ge M$.
\end{assumption}
\noindent
Let $\varphi\in H^1_0(0,T)$ be a fixed nontrivial scalar function and set
\[
U_0 := \{\ p\in L^2(0,T;\R^P): p_\ell(t)=\alpha_\ell\varphi(t)\ \},
\qquad
Y_0 := \{\ G\in L^2(0,T;\R^M): G_i(t)=\beta_i\varphi(t)\ \}.
\]

\begin{proposition}[Approximate right inverse on $Y_0$]
\label{prop:K-right-inverse}
Assume the setting of Theorem~\ref{thm:K-rigorous},
Proposition~\ref{prop:verification-K}, and Assumption~\ref{ass:illumination}.
Then, for $\delta>0$ sufficiently small, the restriction
$\mathcal{K}|_{U_0}:U_0\to Y_0$ is invertible and admits a right-inverse
$\mathcal{K}^\dagger:Y_0\to U_0$ such that
\[
\mathcal{K}\mathcal{K}^\dagger = I_{Y_0},\qquad
\norm{\mathcal{K}^\dagger}_{\mathcal{L}(Y_0,U_0)} \le C,
\]
for some constant $C>0$ independent of $\delta$.
\end{proposition}

\begin{proof}
On $U_0$ and $Y_0$, the operators $\mathcal{T},\mathcal{V},\mathcal{S}$
reduce to finite-dimensional matrices corresponding to the coefficients of
$\varphi(t)$. The leading term of $\mathcal{T}$ is represented by $L_\delta$,
the leading term of $\mathcal{V}$ by an invertible $M\times M$ matrix, and
$\mathcal{S}$ is a diagonal invertible matrix. Thus $\mathcal{K}$ is
represented by $K_\delta:=D M_\delta L_\delta$, where $D$ and $M_\delta$
are invertible and $L_\delta\to L_0$ as $\delta\to0$. Since $L_0$ has rank
$M$ by Assumption~\ref{ass:illumination}, $K_0:=DM_0L_0$ has rank $M$ and
admits a right-inverse $K_0^\dagger$. For $\delta$ small, $K_\delta$ is a
small perturbation of $K_0$, hence admits a right-inverse $K_\delta^\dagger$
with uniformly bounded norm. Identifying $U_0$ and $Y_0$ with $\R^P$ and
$\R^M$ via $\varphi$, we define $\mathcal{K}^\dagger:=K_\delta^\dagger$ and
obtain the stated properties.
\end{proof}

\subsection{Tracking with nanoparticles as actuators}
\label{subsec:tracking-nano}
\noindent
We now combine the abstract closed-loop system and the thermo–plasmonic realization to obtain the tracking result.
\noindent
Let $y_r \in H$ be a fixed reference profile and let $T>0$.
Fix a desired heat input $u_{\mathrm{des}}\in Y_0$ and choose the intensity control $p$ via
$p=\mathcal{K}^\dagger u_{\mathrm{des}}$ as in Proposition~\ref{prop:K-right-inverse}.
Let $y_{Y_0}$ denote the solution of \eqref{eq:5.Y0} driven by $u_{\mathrm{des}}$ and let
$y_{\mathrm{phys}}$ denote the physical temperature generated by the nanoparticle cluster
(i.e.\ by inputs $G$ given by \eqref{eq:K-decomposition}).

\begin{theorem}[Tracking with nanoparticles as actuators (error bounds in $V'$)]
\label{thm:5.7}
Assume the setting of Theorems~\ref{thm:discrete-effective} and~\ref{thm:K-rigorous} and
Proposition~\ref{prop:K-right-inverse}. Let $T>0$ and let $u_{\mathrm{des}}\in Y_0\subset L^2(0,T;\R^M)$.
Choose the intensity control as
\[
  p(t) = \mathcal{K}^\dagger u_{\mathrm{des}}(t)\in U_0.
\]
Let $y_{Y_0}$ be the solution of \eqref{eq:5.Y0}  driven by $u_{\mathrm{des}}$, and let $y_{\mathrm{phys}}$
be the temperature generated by the thermo--plasmonic actuators corresponding to the above choice of $p$.
Then there exist a constant $C_T>0$ and a function $\eta:(0,1]\to(0,\infty)$ with $\eta(\delta)\to 0$
as $\delta\to 0$ such that
\begin{equation}
\label{eq:tracking-estimate}
\| y_{\mathrm{phys}} - y_{Y_0} \|_{C([0,T];V')}^2 + \| y_{\mathrm{phys}} - y_{Y_0} \|_{L^2(0,T;H)}^2
\;\le\; C_T\,\eta(\delta)^2\,\|u_{\mathrm{des}}\|_{L^2(0,T;\R^M)}^2.
\end{equation}
Moreover, for every $t_0\in(0,T]$ there exists $C_{T,t_0}>0$ such that
\begin{equation}
\label{eq:tracking-estimate-H-after-t0}
\sup_{t\in[t_0,T]} \| y_{\mathrm{phys}}(t) - y_{Y_0}(t) \|_{H}
\;\le\; C_{T,t_0}\,\eta(\delta)\,\|u_{\mathrm{des}}\|_{L^2(0,T;\R^M)}.
\end{equation}
\end{theorem}

\begin{proof}
Let $p=\mathcal{K}^\dagger u_{\mathrm{des}}$ and let $G$ be the thermo--plasmonic heat input produced by $p$.
By Theorem~\ref{thm:K-rigorous} and Proposition~\ref{prop:K-right-inverse}, we have
\[
G = \mathcal{K}p + \mathcal{R}^\delta p
  = u_{\mathrm{des}} + \mathcal{R}^\delta p.
\]
Set the input discrepancy $\rho := G-u_{\mathrm{des}}=\mathcal{R}^\delta p$.
Using \eqref{eq:Rdelta-bound} and the boundedness of $\mathcal{K}^\dagger$ on $Y_0$ yields
\[
\|\rho\|_{L^2(0,T;\R^M)}
\le \eta(\delta)\,\|u_{\mathrm{des}}\|_{L^2(0,T;\R^M)},
\]
for some $\eta(\delta)\to 0$ as $\delta\to 0$.
\noindent
Let $e:=y_{\mathrm{phys}}-y_{Y_0}$. Subtracting the equations for $y_{\mathrm{phys}}$ and $y_{Y_0}$ gives
\[
\dot e(t)=A e(t) + B\rho(t),\qquad e(0)=0,
\]
in $D(\mathcal A)'$. Applying Proposition~\ref{prop:wp-Dirac} to this forced equation yields
\[
\|e\|_{C([0,T];V')}^2 + \|e\|_{L^2(0,T;H)}^2
\le C_T\,\|\rho\|_{L^2(0,T;\R^M)}^2,
\]
which implies \eqref{eq:tracking-estimate}. The $H$--estimate on $[t_0,T]$ follows from
\eqref{eq:tracking-estimate} and Corollary~\ref{cor:L2-after-t0}.
\end{proof}

%==================== sec6_error_updates.tex ====================
\section{Robustness: projection and thermo--plasmonic realization errors}
\label{sec:robust-Vprime}
\noindent
The ideal closed-loop error $z_{\mathrm{ideal}}=y_{\mathrm{ideal}}-y_r$ is generated by the feedback
\eqref{eq:KRW-feedback} and solves \eqref{eq:KRW-closed-loop} with the constant forcing $g$ given by
\eqref{eq:feedforward-g}. Let $z_\infty$ be the corresponding equilibrium defined in \eqref{eq:z-infty}.
Then the shifted error $\tilde z_{\mathrm{ideal}}:=z_{\mathrm{ideal}}-z_\infty$ satisfies the homogeneous
stabilized dynamics and hence the exponential estimate in $V'$:
\[
\|\tilde z_{\mathrm{ideal}}(t)\|_{V'} \le e^{-\mu t}\|\tilde z_{\mathrm{ideal}}(0)\|_{V'},\qquad t\ge 0,
\]
for suitable $M,\lambda$ as in Theorem~\ref{thm:KRW-Vprime}.

\subsection{Finite-horizon deviation estimates}
\noindent
Let $z$ be the actual error produced by the thermo--plasmonic actuation scheme.
Assume that the actual forcing differs from the ideal Dirac forcing by a remainder
$r(t)\in D(\mathcal A)'$ on $[0,T]$, i.e.
\[
\dot z(t)=A_0 z(t) -\lambda\sum_{j=1}^M \langle d_j,\mathcal A^{-1}z(t)\rangle\, d_j + g + r(t),
\quad z(0)=z_0\in V'.
\]
Then, by the well-posedness estimate in Proposition~\ref{prop:wp-Dirac} applied to the difference
$e(t):=z(t)-z_{\mathrm{ideal}}(t)$, one obtains
\begin{equation}
\label{eq:err-Vprime}
\|e\|_{C([0,T];V')}^2 + \|e\|_{L^2(0,T;H)}^2
\le C_T\,\|r\|_{L^2(0,T;D(\mathcal A)')}^2.
\end{equation}

\subsection{$L^2$ tracking after a positive time}
\noindent
Fix $t_0\in(0,T]$. Combining \eqref{eq:err-Vprime} with Corollary~\ref{cor:L2-after-t0},
we obtain the $L^2$-in-space deviation estimate on $[t_0,T]$:
\[
\sup_{t\in[t_0,T]}\|z(t)-z_{\mathrm{ideal}}(t)\|_{L^2(\Omega)}
\le C(t_0)\,\|e\|_{L^2((t_0/2,T);H)}
\le C(t_0)\,C_T^{1/2}\,\|r\|_{L^2(0,T;D(\mathcal A)')}.
\]
This is the form in which we will use the robustness estimate in the remainder of the paper.





\section{How Theorem~\ref{thm:5.7} answers the original control problem}
\label{rem:link-original-problem}
\noindent
We divide it into few steps:

\begin{enumerate}
\item The original control question can be formulated as follows: given a desired
temperature profile $y_r$ in $\Omega$, can one design incident electromagnetic
fields acting on a finite collection of thermo--plasmonic nanoparticles so that
the resulting temperature $y_{\mathrm{phys}}(t)$ in the medium tracks $y_r$?

Our analysis proceeds in two steps.

\smallskip
\noindent
(1) \emph{Abstract heat tracking with point actuators.}
In Section~\ref{sec:stab-KRW} we use the feedback law, augmented by a constant feedforward
$u_r$ for non-constant references (see \eqref{eq:ref-finite-modes}--\eqref{eq:feedforward-linear-system}).
The resulting closed loop for the error $z(t)=y(t)-y_r$ is the affine system \eqref{eq:KRW-closed-loop}.
Theorem~\ref{thm:KRW-Vprime} yields exponential stabilization in the $V'$ framework of the shifted error
$z(t)-z_\infty$, and hence convergence of $y_{\mathrm{ideal}}(t)$ towards $y_r+z_\infty$ in $V'$ (and, by smoothing,
in $H$ for any $t\ge t_0>0$).

If $y_r$ is an equilibrium of the free Neumann heat dynamics (equivalently $Ay_r=0$),
then one may take $u_r=0$ and $z_\infty=0$, so that $y_{\mathrm{ideal}}(t)\to y_r$.
In general, for non-equilibrium references $y_r\in X_N$, the above step yields convergence of
$y_{\mathrm{ideal}}(t)$ towards the steady reached profile $y_r+z_\infty$.

To enforce exact tracking of the prescribed low-mode target at steady state, we use the fixed-point
pre-compensation of Subsection~\ref{subsec:fixed-point-bias}. More precisely, we compute a
\emph{commanded reference} $y^\star\in X_N$ by solving the fixed-point equation
\eqref{eq:fixed-point} (see Proposition~\ref{prop:fixed-point}), and we run the same feedback law with reference
$y^\star$ (and corresponding feedforward $u_r(y^\star)=U_Ny^\star$). The resulting closed-loop temperature
converges to the steady reached profile $y_\infty(y^\star)$ defined in \eqref{eq:yinfty-of-ystar}, which satisfies
the exact low-mode matching property $P_Ny_\infty(y^\star)=y_r$ by \eqref{eq:exact-low-mode-tracking}; the
remaining mismatch is a pure tail controlled in Corollary~\ref{cor:fixed-point-tail}.

\smallskip
\noindent
(2) \emph{Realization of the ideal actuators by thermo--plasmonic nanoparticles.}
Sections~\ref{sec:model} and \ref{sec:coupling} show that the cluster of plasmonic nanoparticles, illuminated
by incident fields $p$, generates an effective heat input $G=\mathcal{K}p+\mathcal{R}^\delta p$ in
the background heat equation. Proposition~\ref{prop:K-right-inverse} ensures that on a finite
dimensional subspace $Y_0$ we can invert the leading part of $\mathcal{K}$: for any
desired abstract input $u_{\mathrm{des}}\in Y_0$ there exists an incident field
$p=\mathcal{K}^\dagger u_{\mathrm{des}}$ such that $G=u_{\mathrm{des}}+\mathcal{R}^\delta p$ with
$\|\mathcal{R}^\delta p\|_{L^2(0,T;\mathbb{R}^M)}\le C_T\delta^\mu \|u_{\mathrm{des}}\|$.

\smallskip
Theorem~\ref{thm:5.7} combines these two steps. It compares the ``ideal''
closed-loop temperature $y_{\mathrm{ideal}}$ driven by the abstract control
$u_{\mathrm{des}}$ with the ``physical'' temperature $y_{\mathrm{phys}}$
driven by the thermo--plasmonic actuators corresponding to
$p=\mathcal{K}^\dagger u_{\mathrm{des}}$. The estimate
\[
\sup_{t\in[0,T]} \|y_{\mathrm{phys}}(t)-y_{\mathrm{ideal}}(t)\|_H
  \le C_T\,\eta(\delta)\,\|u_{\mathrm{des}}\|_{L^2(0,T;\mathbb{R}^M)},
\]
with $\eta(\delta)\to 0$ as $\delta\to 0$, shows that the physical temperature
follows the ideal one up to an error that can be made arbitrarily small by
choosing sufficiently small particles (i.e.\ working within the validity
regime of the effective model). Since the closed loop yields exponential decay of the shifted error $z_{\mathrm{ideal}}(t)-z_\infty$ in $V'$ (Theorem~\ref{thm:KRW-Vprime}), combining Theorem~\ref{thm:5.7} and Corollary~\ref{cor:total-error} shows that the nanoparticle--based system drives $y_{\mathrm{phys}}(t)$ towards the corresponding reached steady profile (which equals $y_r$ in the equilibrium case, and satisfies $P_Ny_\infty(y^\star)=y_r$ in the fixed-point compensated case) up to a discrepancy that can be made small by choosing $Y_0$ rich enough and by taking $\delta$ sufficiently small.

In this sense, Theorem~\ref{thm:5.7} provides a rigorous and quantitative
affirmative answer to the original problem of using thermo--plasmonic
nanoparticles as actuators for heat tracking in $\Omega$.

\medskip

\item We now explain how the restriction $u_{\mathrm{des}}\in Y_0$ can be enforced
for a given reference profile $y_r$. Let $\bar y$ denote the reference used in Step~(1): we set $\bar y:=y_r$ in the
equilibrium case, and $\bar y:=y^\star$ in the fixed-point compensated case, where $y^\star$ solves
\eqref{eq:fixed-point}. Let $\bar u_r\in\mathbb{R}^M$ be the corresponding constant feedforward computed as in
\eqref{eq:ref-finite-modes}--\eqref{eq:feedforward-linear-system} (equivalently, $\bar u_r=U_N\bar y$).

Let $u_{\mathrm{ideal}}$ denote the ideal Dirac-actuator input generated by the feedback
(Section~\ref{sec:stab-KRW}), i.e.
\[
  u_{\mathrm{ideal},j}(t)
  := (\bar u_r)_j-\lambda\,\big(\mathcal A^{-1}(y_{\mathrm{ideal}}(t)-\bar y)\big)(x_j),
  \qquad j=1,\dots,M,
\]
where $y_{\mathrm{ideal}}$ is the corresponding closed-loop trajectory.
In general $u_{\mathrm{ideal}}$ lies in $L^2(0,T;\mathbb{R}^M)$, not necessarily in the
finite-dimensional subspace $Y_0$ that can be realized efficiently by the thermo--plasmonic actuators.

A natural way to enforce the constraint $u_{\mathrm{des}}\in Y_0$ is to
project $u_{\mathrm{ideal}}$ onto $Y_0$ and to use only this projected signal
as the ``desired'' input for the thermo--plasmonic realization. More
precisely, we set
\begin{equation}
\label{eq:5.Y0-u-des}
  u_{\mathrm{des}} := P_{Y_0} u_{\mathrm{ideal}},
\end{equation}
where $P_{Y_0}$ denotes the orthogonal projection from
$L^2(0,T;\mathbb{R}^M)$ onto $Y_0$. By construction, $u_{\mathrm{des}}\in Y_0$
and therefore satisfies the hypothesis of Theorem~\ref{thm:5.7}.

Let $y_{Y_0}$ denote the solution of the heat equation driven by
$u_{\mathrm{des}}$:
\begin{equation}
\label{eq:5.Y0}
\begin{cases}
\dot y_{Y_0}(t) = A y_{Y_0}(t) + B u_{\mathrm{des}}(t), & t>0,\\[0.2em]
y_{Y_0}(0) = y_0.
\end{cases}
\end{equation}
The next proposition quantifies the discrepancy between $y_{Y_0}$ and the
ideal closed-loop trajectory $y_{\mathrm{ideal}}$.

\begin{proposition}[Effect of restricting the control to $Y_0$ (in the $V'$ framework)]
\label{prop:projection-error}
Let $T>0$ and let $u_{\mathrm{ideal}}\in L^2(0,T;\R^M)$. Define
$u_{\mathrm{des}}:=P_{Y_0}u_{\mathrm{ideal}}$ as in \eqref{eq:5.Y0-u-des}.
Let $y_{\mathrm{ideal}}$ be the solution of \eqref{eq:heat-DAdual} driven by $u_{\mathrm{ideal}}$
and let $y_{Y_0}$ be the solution of \eqref{eq:5.Y0} driven by $u_{\mathrm{des}}$, with the same initial
condition $y_0\in H\subset V'$. Then there exists $C_T>0$ such that
\begin{equation}
\label{eq:proj-error}
\| y_{Y_0} - y_{\mathrm{ideal}} \|_{C([0,T];V')}^2
+ \| y_{Y_0} - y_{\mathrm{ideal}} \|_{L^2(0,T;H)}^2
\le C_T\,\|(I-P_{Y_0})u_{\mathrm{ideal}}\|_{L^2(0,T;\R^M)}^2.
\end{equation}
Moreover, for every $t_0\in(0,T]$ there exists $C_{T,t_0}>0$ such that
\begin{equation}
\label{eq:proj-error-H-after-t0}
\sup_{t\in[t_0,T]} \| y_{Y_0}(t) - y_{\mathrm{ideal}}(t) \|_{H}
\le C_{T,t_0}\,\|(I-P_{Y_0})u_{\mathrm{ideal}}\|_{L^2(0,T;\R^M)}.
\end{equation}
In particular,
\begin{equation}
\label{eq:best-approx}
\| y_{Y_0} - y_{\mathrm{ideal}} \|_{C([0,T];V')}
\le C_T^{1/2}\inf_{v\in Y_0}\|u_{\mathrm{ideal}}-v\|_{L^2(0,T;\R^M)}.
\end{equation}
\end{proposition}

\begin{proof}
Set $e_{\mathrm{proj}}:=y_{Y_0}-y_{\mathrm{ideal}}$. Then
\[
\dot e_{\mathrm{proj}}(t) = A e_{\mathrm{proj}}(t) + B\bigl(u_{\mathrm{des}}(t)-u_{\mathrm{ideal}}(t)\bigr)
= A e_{\mathrm{proj}}(t) + B\,(I-P_{Y_0})u_{\mathrm{ideal}}(t),
\qquad e_{\mathrm{proj}}(0)=0,
\]
in $D(\mathcal A)'$. Applying Proposition~\ref{prop:wp-Dirac} yields \eqref{eq:proj-error}.
The estimate \eqref{eq:proj-error-H-after-t0} follows from \eqref{eq:proj-error} and
Corollary~\ref{cor:L2-after-t0}. Finally, since $P_{Y_0}$ is the orthogonal projector onto $Y_0$,
\[
\|(I-P_{Y_0})u_{\mathrm{ideal}}\|_{L^2} = \inf_{v\in Y_0}\|u_{\mathrm{ideal}}-v\|_{L^2},
\]
which gives \eqref{eq:best-approx}.
\end{proof}

\medskip
\item Theorem~\ref{thm:5.7} then compares $y_{Y_0}$ with the physical temperature
$y_{\mathrm{phys}}$ generated by the thermo--plasmonic actuators when the
electromagnetic control is chosen as $p = \mathcal{K}^\dagger u_{\mathrm{des}}$.
Combining Proposition~\ref{prop:projection-error} with Theorem~\ref{thm:5.7} yields an
explicit bound on the total tracking error between the ideal closed-loop
trajectory and the physically realized one.

\begin{corollary}[Total tracking error: projection + nanoparticles (in $V'$)]
\label{cor:total-error}
Let $T>0$ and let $u_{\mathrm{ideal}}\in L^2(0,T;\R^M)$ be an ``ideal'' heat input.
Set $u_{\mathrm{des}}:=P_{Y_0}u_{\mathrm{ideal}}$ and let $y_{\mathrm{ideal}}$, $y_{Y_0}$ be as in
Proposition~\ref{prop:projection-error}. Let $y_{\mathrm{phys}}$ be the physical temperature produced by
the thermo--plasmonic actuators when the intensity control $p$ is chosen as in Theorem~\ref{thm:5.7}:
\[
  p = \mathcal{K}^\dagger u_{\mathrm{des}}.
\]
Then there exists $C_T>0$ such that
$$
\|y_{\mathrm{phys}}-y_{\mathrm{ideal}}\|_{C([0,T];V')}^2 + \|y_{\mathrm{phys}}-y_{\mathrm{ideal}}\|_{L^2(0,T;H)}^2
$$
\begin{equation}
\label{eq:total-error}
\le C_T\Bigl(
\|(I-P_{Y_0})u_{\mathrm{ideal}}\|_{L^2(0,T;\R^M)}^2
+ \eta(\delta)^2\,\|u_{\mathrm{des}}\|_{L^2(0,T;\R^M)}^2
\Bigr),
\end{equation}
where $\eta(\delta)\to 0$ is as in Theorem~\ref{thm:5.7}. Moreover, for every $t_0\in(0,T]$,
\[
\sup_{t\in[t_0,T]}\|y_{\mathrm{phys}}(t)-y_{\mathrm{ideal}}(t)\|_{H}
\le C_{T,t_0}\Bigl(
\|(I-P_{Y_0})u_{\mathrm{ideal}}\|_{L^2(0,T;\R^M)}
+ \eta(\delta)\,\|u_{\mathrm{des}}\|_{L^2(0,T;\R^M)}
\Bigr),
\]
for some constant $C_{T,t_0}>0$.
\end{corollary}

\begin{proof}
We decompose
\[
y_{\mathrm{phys}}-y_{\mathrm{ideal}}=(y_{\mathrm{phys}}-y_{Y_0})+(y_{Y_0}-y_{\mathrm{ideal}}).
\]
Applying Theorem~\ref{thm:5.7} to $y_{\mathrm{phys}}-y_{Y_0}$ and Proposition~\ref{prop:projection-error}
to $y_{Y_0}-y_{\mathrm{ideal}}$, and then using the triangle inequality, yields \eqref{eq:total-error}.
The $H$--estimate on $[t_0,T]$ follows similarly from \eqref{eq:tracking-estimate-H-after-t0} and
\eqref{eq:proj-error-H-after-t0}.
\end{proof}
\end{enumerate}

\medskip

\noindent
Proposition~\ref{prop:projection-error} measures
the error introduced by restricting the abstract feedback design to controls
lying in the finite-dimensional subspace $Y_0$, while Theorem~\ref{thm:5.7} measures
the additional error incurred when realizing such controls by thermo--plasmonic
nanoparticles. Corollary~\ref{cor:total-error} shows that the total tracking
error is the sum of these two contributions: a purely ``approximation'' term,
given by the best $L^2$--approximation of $u_{\mathrm{ideal}}$ in $Y_0$, and
a ``physical realization'' term of order $\eta(\delta)$ coming from the
effective thermo--plasmonic model.







%========================================================
\section{Algorithmic design}
\label{sec:algorithm}
%========================================================
\noindent
In this section, we summarize how the abstract results of Sections~\ref{sec:feedback}--\ref{sec:coupling} can be turned into a constructive procedure for heat tracking using thermo--plasmonic nanoparticles.
\newline
The feedback design in Section~\ref{sec:feedback} produces an \emph{ideal} closed-loop
trajectory $y_{\mathrm{ideal}}$ for the abstract point-actuator model, driven
by an implicit control signal $u_{\mathrm{ideal}}$. The thermo--plasmonic
analysis in Section~\ref{sec:coupling} shows how to realize a finite-dimensional approximation
$u_{\mathrm{des}}\in Y_0$ of this signal by suitable EM illuminations
$p\in U_0$, with a small error controlled by the particle size $\delta$.
\newline
We now discuss how each ingredient can be constructed in practice and how
they combine into an algorithm.


\subsection{Implementable feedback law}
\label{subsec:alg-KRW}
\noindent
Let $y_r\in X_N$ satisfy \eqref{eq:ref-finite-modes}. For equilibrium references (i.e.\ $A_0y_r=0$), we set
$\bar y:=y_r$ and $u_r:=0$. For general (non-equilibrium) references, we incorporate the fixed-point
pre-compensation of Subsection~\ref{subsec:fixed-point-bias}: we compute the \emph{commanded reference}
$y^\star\in X_N$ as the solution of \eqref{eq:fixed-point} (see Proposition~\ref{prop:fixed-point}) and set
$\bar y:=y^\star$. In both cases, we compute a constant feedforward vector $u_r\in\R^M$ associated with
$\bar y$ by solving \eqref{eq:feedforward-linear-system} (equivalently $u_r=U_N\bar y$ in the notation of Section~\ref{sec:stab-KRW}).
\newline
We use the feedback
\[
u_j(t) = (u_r)_j-\lambda\,\big(\mathcal A^{-1}(y(t)-\bar y)\big)(x_j),
\qquad \mathcal A := I-\kappa_m\Delta,
\]
which is equivalent to \eqref{eq:KRW-feedback} (with reference $\bar y$).
In the fixed-point compensated case, the reached steady profile $y_\infty(\bar y)$ satisfies
$P_Ny_\infty(\bar y)=y_r$ by \eqref{eq:exact-low-mode-tracking}, and the remaining mismatch is a pure tail
controlled by Corollary~\ref{cor:fixed-point-tail}.

\paragraph{Offline choice of actuator points and $\lambda$.}
Choose actuator locations $x_1,\dots,x_M$ so that Assumption~\ref{assump:KRW-act} holds.
In rectangular domains, \cite{KunRodWal24-cocv} provides explicit constructive placements.
Then, for a desired decay rate $\mu>0$, choose $M$ and $\lambda>0$ according to Theorem~\ref{thm:KRW-Vprime}.
(Practically, $\lambda$ can be tuned numerically by monitoring the decay of $\|y(t)-\bar y\|_{V'}$.)

\paragraph{Online step (computing $u(t)$).}
At each time $t$:
\begin{enumerate}
\item Form the error state $z(t)=y(t)-\bar y$.
\item Solve the elliptic problem for $w(t)=\mathcal A^{-1}z(t)$:
\[
w(t) - \kappa_m\Delta w(t) = z(t)\ \text{in }\Omega,\qquad \partial_\nu w(t)=0\ \text{on }\partial\Omega.
\]
\item Set $u_j(t)=(u_r)_j-\lambda\,w(t,x_j)$ for $j=1,\dots,M$.
\end{enumerate}

\paragraph{Remark on sensors.}
The above requires knowledge of $z(t)$ in $\Omega$ to solve the elliptic problem.
If only boundary (or locally averaged) measurements are available, one needs a state estimation step; this lies beyond
the present scope and can be incorporated by standard observer constructions once the model is fixed.


\subsection{Identification of the thermo--plasmonic map $\mathcal{K}$}
\label{subsec:alg-Ktp}
%--------------------------------------------------------
\noindent
The thermo--plasmonic analysis of Section~\ref{sec:coupling} shows that, in the effective
regime covered by Theorem~\ref{thm:K-rigorous}, the cluster of nanoparticles acts on the
background heat equation as an effective finite number of point sources:
\begin{equation}
\label{eq:alg-G-Ktp}
G = \mathcal{K} p + R_\delta p,
\end{equation}
where $p$ is the intensity control (incident fields), $G$ is the resulting heat
input in the background heat equation, $\mathcal{K}$ is a bounded linear
operator (explicitly decomposed as $SVT$ in Theorem~\ref{thm:K-rigorous}), and $R_\delta$ is
a small remainder of order $\mathcal{O}(\delta^\mu)$ in suitable norms.
\noindent
We restrict our attention to finite-dimensional subspaces
$U_0\subset U_{\mathrm{em}}$ and $Y_0\subset L^2(0,T;\R^M)$ as in
Section~\ref{sec:map-EM-Heat}, so that $\mathcal{K}$ is represented by a matrix
$K_0:U_0\to Y_0$. Theorem~\ref{thm:K-rigorous} and Proposition~\ref{prop:K-right-inverse} guarantee that, under the
assumptions on the geometry, material parameters and illuminations,
$K_0$ has full rank and admits a bounded right inverse $K_0^\dagger$.
\noindent
In practice, $\mathcal{K}$ is most conveniently obtained by as follows:

\begin{enumerate}
  \item Choose a basis $\{\psi^{(1)},\dots,\psi^{(P)}\}$ of $U_0$ (e.g.\ a
  finite set of time profiles and spatial illuminations) and identify
  $U_0\simeq\R^P$ via
  \[
    p(t) = \sum_{\ell=1}^P \alpha_\ell \psi^{(\ell)}(t)
    \quad\Longleftrightarrow\quad
    \alpha = (\alpha_1,\dots,\alpha_P)^\top\in\R^P.
  \]

  \item Fix a basis of $Y_0$ (for example $Y_0=\{\beta\varphi(t):\beta\in\R^M\}$
  with a given time profile $\varphi$) and identify $Y_0\simeq\R^M$ in the
  same way.

  \item For each basis illumination $\psi^{(\ell)}$, solve (numerically) the
  coupled Maxwell--heat model in the effective regime and extract the
  resulting heat input $G^{(\ell)}\in Y_0$ entering the background heat
  equation. In the matrix representation, this corresponds to
  \[
    K_0 e_\ell = g^{(\ell)},\qquad
    g^{(\ell)}\in\R^M,
  \]
  where $e_\ell$ is the $\ell$th canonical basis vector in $\R^P$ and
  $g^{(\ell)}$ is the coordinate vector of $G^{(\ell)}$ in $Y_0$.

  \item Assemble the matrix $K_0\in\R^{M\times P}$ whose columns are
  $g^{(1)},\dots,g^{(P)}$. The structural analysis in Section~\ref{sec:coupling} ensures that,
  for a suitable choice of illuminations and small $\delta$, $K_0$ has full
  rank $M$ and $\|K_0^\dagger\|$ is uniformly bounded.

  \item Compute a right inverse $K_0^\dagger:Y_0\to U_0$, for example via an
  SVD-based pseudo-inverse. This is the constructive version of the operator
  $K^\dagger$ used in Theorem~\ref{thm:5.7}.
\end{enumerate}
Once $K_0$ and $K_0^\dagger$ are known, the finite-dimensional relation
\eqref{eq:alg-G-Ktp} can be used pointwise in time as
\[
  G(t) \approx K_0 p(t),\qquad
  p(t) \approx K_0^\dagger G(t),
\]
with an error of order $\mathcal{O}(\delta^\mu)$ quantified in
Theorem~\ref{thm:5.7}.

%--------------------------------------------------------
\subsection{Control loop with projection onto $Y_0$}
\label{subsec:alg-loop}
%--------------------------------------------------------
\noindent
Let $y_r$ satisfy \eqref{eq:ref-finite-modes}. Define the reference $\bar y$ used in the feedback as follows:
set $\bar y:=y_r$ in the equilibrium case ($A_0y_r=0$), and $\bar y:=y^\star$ in the non-equilibrium case,
where $y^\star\in X_N$ solves the fixed-point equation \eqref{eq:fixed-point}.
Compute a constant feedforward vector $u_r\in\R^M$ associated with $\bar y$ by solving
\eqref{eq:feedforward-linear-system} (equivalently, $u_r=U_N\bar y$).
\noindent
Let $u_{\mathrm{ideal}}$ be the \emph{ideal} Dirac-actuator input generated by the feedback law
(Section~\ref{sec:stab-KRW}):
\begin{equation}
\label{eq:alg-u-ideal}
u_{\mathrm{ideal},j}(t) := (u_r)_j -\lambda\,\big(\mathcal A^{-1}(y(t)-\bar y)\big)(x_j),
\qquad j=1,\dots,M,\quad \mathcal A:=I-\kappa_m\Delta.
\end{equation}
In general $u_{\mathrm{ideal}}\notin Y_0$; we therefore project it onto $Y_0$ as in \eqref{eq:5.Y0-u-des}:
\begin{equation}
\label{eq:alg-u-des}
u_{\mathrm{des}} := P_{Y_0} u_{\mathrm{ideal}} \in Y_0.
\end{equation}
The auxiliary trajectory $y_{Y_0}$ driven by $u_{\mathrm{des}}$ is defined by \eqref{eq:5.Y0}.
Proposition~\ref{prop:projection-error} quantifies the deviation between $y_{Y_0}$ and the ideal trajectory driven by $u_{\mathrm{ideal}}$ in the $V'$ framework (and in $H$ after any $t_0>0$).
\noindent
To realize $u_{\mathrm{des}}$ physically, we choose the intensity control $p\in U_0$ via
\begin{equation}
\label{eq:alg-p-Kdagger}
p := K_0^\dagger u_{\mathrm{des}},
\end{equation}
so that $K_0 p$ approximates $u_{\mathrm{des}}$ in the sense of \eqref{eq:alg-G-Ktp}.
Theorem~\ref{thm:5.7} then yields bounds on $y_{\mathrm{phys}}-y_{Y_0}$ in $V'$ (and in $H$ after $t_0>0$).
Combining both contributions gives Corollary~\ref{cor:total-error}.

\medskip

\noindent\textbf{Summary of the implementable loop.}
\begin{enumerate}
\item \textbf{Choose actuator points and gain.}
Select $x_1,\dots,x_M$ so that Assumption~\ref{assump:KRW-act} holds and choose $\lambda>0$
according to Theorem~\ref{thm:KRW-Vprime} for the desired decay rate in $V'$.

\item \textbf{(If needed) fixed-point pre-compensation on $X_N$.}
If $y_r\in X_N$ is not an equilibrium of the free Neumann heat dynamics, solve \eqref{eq:fixed-point} for
$y^\star$ and set $\bar y:=y^\star$; otherwise set $\bar y:=y_r$.
Compute the corresponding feedforward $u_r$ by solving \eqref{eq:feedforward-linear-system}
(equivalently $u_r=U_N\bar y$).

\item \textbf{Identify $K_0$ and $K_0^\dagger$.}
Fix finite-dimensional spaces $U_0$, $Y_0$ and perform the calibration described in
Subsection~\ref{subsec:alg-Ktp} to obtain $K_0$ and a right inverse $K_0^\dagger$.

\item \textbf{At each time $t$ (feedback + projection + EM actuation).}
\begin{enumerate}
\item Form the error state $z(t)=y(t)-\bar y$.
\item Solve $w(t)=\mathcal A^{-1}z(t)$, i.e.\ $w-\kappa_m\Delta w=z$ in $\Omega$ with $\partial_\nu w=0$ on $\partial\Omega$.
\item Set $u_{\mathrm{ideal},j}(t)=(u_r)_j-\lambda\,w(t,x_j)$.
\item Project $u_{\mathrm{des}}(t)=P_{Y_0}u_{\mathrm{ideal}}(t)$.
\item Compute $p(t)=K_0^\dagger u_{\mathrm{des}}(t)$ and apply the corresponding EM illumination during the next time step.
\end{enumerate}
\end{enumerate}


\section{Appendix}\label{appendix}
\subsection{Spectral assumptions}
\begin{proposition}[Genericity of Assumption~\ref{ass:rank-4}]
\label{prop:generic-rank}
Fix $M\in\mathbb{N}$ and consider the map
\[
F:(\Omega)^M \to \mathbb{R},\qquad
F(x_1,\dots,x_M) := \det\big(\phi_k(x_j)\big)_{j,k=1}^M,
\]
where $\{\phi_k\}_{k\ge1}$ is the eigenbasis of $-A_0$ introduced above.
Then:
\begin{enumerate}
\item[(a)] $F$ is a real-analytic function on $(\Omega)^M$.
\item[(b)] $F$ is not identically zero.
\item[(c)] The zero set 
\[
\mathcal{Z}_M := \{(x_1,\dots,x_M)\in\Omega^M : F(x_1,\dots,x_M)=0\}
\]
has Lebesgue measure zero in $\Omega^M$.
\end{enumerate}
In particular, Assumption~\ref{ass:rank-4} holds for almost every choice of
the actuator positions $(x_1,\dots,x_M)\in\Omega^M$ in the sense of
Lebesgue measure. Then Assumption ~\ref{ass:rank-4} is a full row rank condition for the 
$N\times M$ matrix of the first $
N$ modes, with $N \leq M$.
\end{proposition}
\noindent This provides a genericity mechanism for point actuator/sensor placement. We include the proof for completeness, since we require the specific full-row-rank condition in Assumption~\ref{ass:rank-4}(i) on the first $N$ modes.
 


%This is a standard genericity mechanism for point actuator/sensor placement. We include the proof for completeness since we require the specific rank condition on the first $N$ modes.
\begin{proof}
(a) Each eigenfunction $\phi_k$ is real-analytic in $\Omega$ (since $A$ is
a second-order elliptic operator with smooth coefficients (i.e. constant here) and Neumann boundary
conditions). The entries $\phi_k(x_j)$ of the matrix thus depend
real-analytically on $(x_1,\dots,x_M)$, and the determinant is a polynomial
expression in these entries; hence $F$ is real-analytic.

(b) Suppose, for contradiction, that $F\equiv 0$ on $\Omega^M$. Then for
every $M$-tuple $(x_1,\dots,x_M)\in\Omega^M$ the $M\times M$ matrix
$(\phi_k(x_j))_{j,k=1}^M$ is singular, i.e.\ its rows
\[
r(x_j) := \big(\phi_1(x_j),\dots,\phi_M(x_j)\big)\in\mathbb{R}^M
\]
are linearly dependent. This implies that the set
\[
\mathcal{R} := \{\, r(x) : x\in\Omega \,\} \subset \mathbb{R}^M
\]
is contained in a proper linear subspace of $\mathbb{R}^M$. Hence there
exists a nonzero vector $c=(c_1,\dots,c_M)\in\mathbb{R}^M$ such that
\[
c\cdot r(x) = \sum_{k=1}^M c_k \phi_k(x) = 0,\qquad x\in\Omega.
\]
Thus $\sum_{k=1}^M c_k \phi_k \equiv 0$ in $\Omega$. Since
$\{\phi_1,\dots,\phi_M\}$ are linearly independent in $H$, this forces
$c_k=0$ for all $k$, a contradiction. Therefore $F$ is not identically zero.

(c) The zero set of a nontrivial real-analytic
function has Lebesgue measure zero. Since $F$ is real-analytic and not
identically zero, the set $\mathcal{Z}_M$ has measure zero in $\Omega^M$.
\end{proof}

\begin{remark}[Practical interpretation]
\label{rem:practical-rank}
Proposition~\ref{prop:generic-rank} shows that, for fixed $M$ and domain
$\Omega$, Assumption~\ref{ass:rank-4} is a generic property of the actuator
positions $(x_1,\dots,x_M)$: it fails only on an exceptional set
$\mathcal{Z}_M$ of Lebesgue measure zero. Equivalently, if the points
$x_j$ are chosen ``at random'' in $\Omega$ (according to any absolutely
continuous probability distribution), then the matrix
$\mathcal{B}_{jk}=\phi_k(x_j)$ is invertible with probability one.
\newline
In particular, the rank condition can always be enforced in practice by
perturbing a given configuration of points slightly (as long as the points
remain in the interior of $\Omega$ and do not approach the boundary or
coincide). In simple geometries, such as intervals or rectangles where the
eigenfunctions are explicitly known (e.g.\ trigonometric functions), it is
enough to ensure that:
\begin{itemize}
\item the points $x_j$ do not lie on nodal sets of any of the eigenfunctions
$\phi_k$ with $k\le M$, and
\item the points are not arranged in a way that respects the symmetries of
these nodal sets.
\end{itemize}
For example, on an interval $\Omega=(0,L)$ with Neumann boundary conditions,
the eigenfunctions are $\phi_1\equiv \mathrm{const}$ and
$\phi_k(x)=\cos((k-1)\pi x/L)$ for $k\ge2$. Choosing $M$ distinct points
$x_j\in(0,L)$ that avoid the zeros of these cosine functions and are not
symmetrically paired with respect to the midpoint generically yields an
invertible matrix $\mathcal{B}$.
\end{remark}
\subsubsection*{A concrete condition in 1D (Neumann interval)}
\noindent
We first treat the case of a one-dimensional interval with Neumann boundary
conditions, for which the eigenfunctions are explicitly known.

\begin{lemma}[Invertibility on a 1D Neumann interval via DCT nodes]
\label{lem:rank-1D}
Let $\Omega=(0,L)$ and let $A=\kappa_m\partial_{xx}$ with Neumann boundary
conditions. An orthonormal basis of eigenfunctions of $-A_0$ is given by
\[
\phi_1(x) := \frac{1}{\sqrt{L}},\qquad
\phi_k(x) := \sqrt{\frac{2}{L}}\,\cos\!\Big(\frac{(k-1)\pi x}{L}\Big),
\quad k\ge2.
\]
Fix $M\in\mathbb{N}$ and define actuator positions
\begin{equation}
\label{eq:1D-nodes}
x_j := \frac{(2j-1)L}{2M},\qquad j=1,\dots,M.
\end{equation}
Let $\mathcal{B}\in\mathbb{R}^{M\times M}$ be the matrix
\[
\mathcal{B}_{jk} := \phi_k(x_j),\qquad j,k=1,\dots,M.
\]
Then $\mathcal{B}$ is invertible.
\end{lemma}

\begin{proof}
By the explicit formulas for $\phi_k$ and \eqref{eq:1D-nodes}, we have
\[
\phi_1(x_j) = \frac{1}{\sqrt{L}},\qquad
\phi_k(x_j) = \sqrt{\frac{2}{L}}\,
\cos\!\Big(\frac{(k-1)\pi(2j-1)}{2M}\Big),\quad k\ge2.
\]
Factor out the (nonzero) normalizing constants row- and column-wise. Up to
multiplication by invertible diagonal matrices on the left and right,
$\mathcal{B}$ is equivalent to the matrix $C\in\mathbb{R}^{M\times M}$ with
entries
\[
C_{jk} := \cos\!\Big(\frac{(k-1)\pi(2j-1)}{2M}\Big),
\qquad j,k=1,\dots,M.
\]
The matrix $C$ is (up to normalization) the discrete cosine transform of type
II (DCT-II), known to be orthogonal (after an appropriate scaling of the first
column), see \cite{StrangBanks} for instance. In particular, $\det C\neq0$. Since multiplication by invertible
diagonal matrices does not change rank, we conclude that $\mathcal{B}$ is
invertible.
\end{proof}

%\begin{remark}[Choice of $H_M$ in 1D]
%\label{rem:HM-1D}
%In Lemma~\ref{lem:rank-1D} we took $H_M=\mathrm{span}\{\phi_1,\dots,\phi_M\}$,
%i.e.\ the first $M$ Neumann eigenfunctions ordered by increasing eigenvalue,
%and actuator positions \eqref{eq:1D-nodes}. Lemma~\ref{lem:rank-1D} shows that
%in 1D. The nodes \eqref{eq:1D-nodes} are simply the midpoints of the $M$
%equal subintervals of $(0,L)$, matching the standard DCT-II sampling grid.
%\end{remark}

\subsubsection*{A concrete condition in 3D for a Neumann cube}
\noindent
We now give an explicit sufficient condition for a three-dimensional cube,
using tensor products of the 1D construction.

\begin{lemma}[Invertibility on a 3D Neumann cube via tensor-product DCT grid]
\label{lem:rank-3D}
Let $\Omega=(0,L_1)\times(0,L_2)\times(0,L_3)$ and let
$A=\kappa_m\Delta$ with Neumann boundary conditions on $\partial\Omega$.
An orthonormal basis of eigenfunctions of $-A_0$ is given (up to
normalization) by
\[
\phi_{\mathbf{n}}(x)
= c_{\mathbf{n}}\,
\cos\!\Big(\frac{n_1\pi x_1}{L_1}\Big)
\cos\!\Big(\frac{n_2\pi x_2}{L_2}\Big)
\cos\!\Big(\frac{n_3\pi x_3}{L_3}\Big),
\]
where $\mathbf{n}=(n_1,n_2,n_3)\in\mathbb{N}_0^3$ and $c_{\mathbf{n}}>0$ are
normalizing constants.
\newline
Fix integers $N_1,N_2,N_3\ge0$ and set
\[
I := \{\,\mathbf{n}=(n_1,n_2,n_3)\in\mathbb{N}_0^3 : 0\le n_\ell\le N_\ell,
\ \ell=1,2,3 \,\}.
\]
Let
\[
H_M := \mathrm{span}\{\phi_{\mathbf{n}} : \mathbf{n}\in I\},
\quad M:=(N_1+1)(N_2+1)(N_3+1).
\]
Index $H_M$ by an enumeration $\{\phi_k\}_{k=1}^M$ of $\{\phi_{\mathbf{n}} :
\mathbf{n}\in I\}$ (the order is irrelevant for rank). For each direction
$\ell=1,2,3$ define
\[
x^{(\ell)}_{j_\ell} := \frac{(2j_\ell-1)L_\ell}{2(N_\ell+1)},
\qquad j_\ell=1,\dots,N_\ell+1,
\]
and form the tensor-product grid
\[
x_{\mathbf{j}}
:= \big(x^{(1)}_{j_1},x^{(2)}_{j_2},x^{(3)}_{j_3}\big),
\qquad
\mathbf{j}=(j_1,j_2,j_3)\in J,
\]
where $J:=\{1,\dots,N_1+1\}\times\{1,\dots,N_2+1\}\times\{1,\dots,N_3+1\}$
and $|J|=M$.
\newline
Let $\mathcal{B}\in\mathbb{R}^{M\times M}$ be the matrix
\[
\mathcal{B}_{\mathbf{j},\mathbf{n}} := \phi_{\mathbf{n}}(x_{\mathbf{j}}),
\qquad \mathbf{j}\in J,\ \mathbf{n}\in I,
\]
and identify $\mathcal{B}$ with an $M\times M$ matrix after any fixed
enumeration of $J$ and $I$. Then $\mathcal{B}$ is invertible.
\end{lemma}

\begin{proof}
By the tensor-product structure of the eigenfunctions and the grid, we have
\[
\phi_{\mathbf{n}}(x_{\mathbf{j}})
= c_{\mathbf{n}}\,
\prod_{\ell=1}^3 \cos\!\Big(\frac{n_\ell\pi x^{(\ell)}_{j_\ell}}{L_\ell}\Big).
\]
Define, for each $\ell=1,2,3$, the 1D matrices
\[
C^{(\ell)}_{j_\ell,n_\ell}
:= \cos\!\Big(\frac{n_\ell\pi x^{(\ell)}_{j_\ell}}{L_\ell}\Big),
\qquad j_\ell=1,\dots,N_\ell+1,\quad n_\ell=0,\dots,N_\ell.
\]
By the same argument as in Lemma~\ref{lem:rank-1D} (with $M$ replaced by
$N_\ell+1$), each $C^{(\ell)}$ is, up to normalization, a DCT-II matrix and
hence invertible.
\newline
Up to multiplication by nonzero diagonal matrices corresponding to the
normalization constants $c_{\mathbf{n}}$, the full matrix $\mathcal{B}$ can
be written as a (permuted) Kronecker product
\[
\mathcal{B} = D \cdot\big(C^{(1)}\otimes C^{(2)}\otimes C^{(3)}\big),
\]
where $D$ is a diagonal matrix with strictly positive entries and
$\otimes$ denotes the Kronecker product. Since each $C^{(\ell)}$ is
invertible and the Kronecker product of invertible matrices is invertible,
$C^{(1)}\otimes C^{(2)}\otimes C^{(3)}$ is invertible. Multiplication by
the diagonal matrix $D$ does not affect invertibility. Therefore
$\mathcal{B}$ is invertible. In particular, any selection of $N\le M$ rows yields an $N\times M$ matrix with full row rank~$N$.
\end{proof}









\subsection{Verification of Assumption~\ref{assump:KRW-act} on $C^{1,1}$ domains}
\label{subsec:app-KRW-act-C11}
%--------------------------------------------------------
\noindent
This subsection provides a self-contained verification of the actuator hypothesis
(Assumption~\ref{assump:KRW-act}) for Dirac actuators on smooth domains.
The only part of Assumption~\ref{assump:KRW-act} that is geometric is the divergence
$\xi_M^+\to+\infty$ of the nullspace coercivity constant.
In \cite{KunRodWal24-cocv} the authors verify this in rectangles/cubes
by a self-similar decomposition argument. Here we replace that step by a standard
finite-element interpolation estimate on shape-regular quasi-uniform meshes, which works on general
$C^{1,1}$ domains (hence extending the class of box/polyhedral domains  in \cite{KunRodWal24-cocv} to a class of domains admitting a suitable meshing property; see~\eqref{eq:app-shapereg}).

\medskip

\noindent
Throughout, $\Omega\subset\R^d$ is bounded of class $C^{1,1}$ with $d\in\{1,2,3\}$, $A_0=\kappa_m\Delta$ with
homogeneous Neumann boundary condition, and $\mathcal A=I-A_0=I-\kappa_m\Delta$ as in the main text.
We recall that $D(\mathcal A)=\{w\in H^2(\Omega):\partial_\nu w=0\text{ on }\partial\Omega\}$ and
$V=H^1(\Omega)$.

\begin{remark}[Actuator points on $\overline\Omega$] 
\label{rem:app-KRW-closure}
In the main text we fix actuator locations $x_j\in\Omega$ for simplicity. The functional-analytic framework used
throughout only requires $\delta_x\in D(\mathcal A)'$, which holds for every $x\in\overline\Omega$ because
$D(\mathcal A)\hookrightarrow C(\overline\Omega)$ (see Proposition~\ref{prop:B-bounded-DA-dual}).
Therefore, in the verification below we may select actuator points among the vertices of a triangulation, which may
include boundary vertices. This does not affect Assumption~\ref{assump:KRW-act} nor the stabilization statement.
\end{remark}

\begin{lemma}[Elliptic graph norm equivalence]
\label{lem:app-graphnorm-H2}
There exist constants $c_\mathcal{A},C_\mathcal{A}>0$ depending only on $\Omega$ and $\kappa_m$ such that
\begin{equation}
\label{eq:app-graphnorm-H2}
c_\mathcal{A}\,\|w\|_{H^2(\Omega)} \le \|w\|_{D(\mathcal A)} \le C_\mathcal{A}\,\|w\|_{H^2(\Omega)}
\qquad \forall w\in D(\mathcal A).
\end{equation}
In particular, point evaluation $w\mapsto w(x)$ is continuous on $D(\mathcal A)$ for every $x\in\overline\Omega$.
\end{lemma}

\begin{proof}
The upper bound in \eqref{eq:app-graphnorm-H2} follows from $\mathcal A w=w-\kappa_m\Delta w$ and the estimate
$\|w\|_{L^2}+\|\Delta w\|_{L^2}\lesssim\|w\|_{H^2}$. For the reverse inequality, let $f:=\mathcal A w\in L^2(\Omega)$.
Then $w$ solves the Neumann problem $(I-\kappa_m\Delta)w=f$.
On bounded $C^{1,1}$ domains, $H^2$-regularity holds for this uniformly elliptic operator with Neumann boundary condition,
so $\|w\|_{H^2}\le C\|f\|_{L^2}=C\|w\|_{D(\mathcal A)}$.
The continuity of point evaluation follows from the embedding $H^2(\Omega)\hookrightarrow C(\overline\Omega)$ for $d\le3$.
\end{proof}

\medskip
\noindent\textbf{Choice of actuator families.}
Let $\{\mathcal T_h\}_{h\downarrow 0}$ be a family of shape-regular, quasi-uniform simplicial triangulations of $\Omega$.
Denote by $\mathcal N_h=\{x_{h,1},\dots,x_{h,M(h)}\}\subset\overline\Omega$ the set of all mesh vertices.
For each $h$ (equivalently, for each $M:=M(h)$), define Dirac actuators
\begin{equation}
\label{eq:app-actuators-dirac}
d_{M,j}:=\delta_{x_{h,j}}\in D(\mathcal A)',\qquad j=1,\dots,M.
\end{equation}
To produce biorthogonal test functions in $D(\mathcal A)$, we use Riesz representatives.

\begin{lemma}[Biorthogonal family in $D(\mathcal A)$]
\label{lem:app-biorthogonal}
For each $M$ and each $j\in\{1,\dots,M\}$ there exists $\Psi_{M,j}\in D(\mathcal A)$ such that
\begin{equation}
\label{eq:app-biorthogonality}
\langle d_{M,i},\Psi_{M,j}\rangle = \Psi_{M,j}(x_{h,i}) = \delta_{ij},\qquad i,j=1,\dots,M.
\end{equation}
Consequently, $\dim\mathrm{span}\{d_{M,j}\}_{j=1}^M=\dim\mathrm{span}\{\Psi_{M,j}\}_{j=1}^M=M$.
\end{lemma}

\begin{proof}
By Lemma~\ref{lem:app-graphnorm-H2}, each evaluation functional $d_{M,j}$ is continuous on the Hilbert space
$D(\mathcal A)$ equipped with the inner product $(u,v)_{D(\mathcal A)}:=(\mathcal A u,\mathcal A v)_H$.
Hence, by the Riesz representation theorem, there exists a unique $r_{M,j}\in D(\mathcal A)$ such that
$\langle d_{M,j},w\rangle=(w,r_{M,j})_{D(\mathcal A)}$ for all $w\in D(\mathcal A)$.
Define the Gram matrix $G_M\in\R^{M\times M}$ by
\[
(G_M)_{ij}:=\langle d_{M,i},r_{M,j}\rangle = (r_{M,i},r_{M,j})_{D(\mathcal A)}.
\]
For any $c\in\R^M$, we have $c^\top G_M c=\big\|\sum_{j=1}^M c_j r_{M,j}\big\|_{D(\mathcal A)}^2$, so $G_M$ is symmetric
positive definite and thus invertible. Now set
\[
\Psi_{M,j} := \sum_{\ell=1}^M (G_M^{-1})_{j\ell}\, r_{M,\ell}\in D(\mathcal A).
\]
Then, for every $i$,
\[
\langle d_{M,i},\Psi_{M,j}\rangle
=\sum_{\ell=1}^M (G_M^{-1})_{j\ell}\,\langle d_{M,i},r_{M,\ell}\rangle
=\sum_{\ell=1}^M (G_M^{-1})_{j\ell}\,(G_M)_{i\ell} = \delta_{ij},
\]
which is \eqref{eq:app-biorthogonality}. The dimension claims follow immediately from biorthogonality.
\end{proof}

\medskip
\noindent\textbf{Divergence of the nullspace coercivity constant.}
Let $V_h\subset H^1(\Omega)$ be the standard piecewise-affine ($P_1$) finite element space associated with $\mathcal T_h$,
and let $I_h:H^2(\Omega)\to V_h$ be the nodal Lagrange interpolant.

\begin{lemma}[Nodal interpolation estimate for polyhedral domains]
\label{lem:app-interp}
There exists a constant $C_I>0$, independent of $h$, such that for all $w\in H^2(\Omega)$,
\begin{equation}
\label{eq:app-interp}
\|w-I_h w\|_{H^1(\Omega)} \le C_I\, h\,\|w\|_{H^2(\Omega)}.
\end{equation}
\end{lemma}

\begin{proof}
This is the classical $P_1$ interpolation estimate on shape-regular quasi-uniform meshes; see Theorem $4.4.20$ in \cite{BrennerScottFEM} for polyhedral domains.
\end{proof}

%%%%%%%%%%%%%%%%%%%%%%%%%%%%%%%%%%%%%%%%%%%%%%%%%%%
%%%%%%%%%%%%%%%%%%%%%%%%%%%%%%%%%%%%%%%%%%%%%%%%%%%
% We use curved instead of affine FEM-spaces 
%%%%%%%%%%%%%%%%%%%%%%%%%%%%%%%%%%%%%%%%%%%%%%%%%
%%%%%%%%%%%%%%%%%%%%%%%%%%%%%%%%%%%%%%%%%%%%%%%%%


\medskip
 \noindent To deal with more general domains as polyhedral ones, we take $\Omega\subset\R^3$ to be a bounded $C^{1,1}$ domain and work with a boundary-fitted isoparametric mesh family (see, e.g., \cite{Lenoir1986,CiarletRaviart1972, AzouaniTiti14,LunasinTiti17,EggerFritzKunRod25})
$\{\mathcal T_h\}_{h\downarrow 0}$: each element $K\in\mathcal T_h$ is the image of a fixed reference simplex
$\widehat K$ under a bijective map $F_K\in W^{2,\infty}(\widehat K;\R^3)$.
We assume uniform (isoparametric) shape-regularity and quasi-uniformity: there exist constants $\gamma,c_0>0$,
independent of $h$ and $K$, such that
\begin{equation}\label{eq:app-shapereg}
\|DF_K\|_{L^\infty(\widehat K)}\le \gamma h_K,\qquad
\|(DF_K)^{-1}\|_{L^\infty(\widehat K)}\le \gamma h_K^{-1},\qquad
\|D^2F_K\|_{L^\infty(\widehat K)}\le \gamma h_K,
\end{equation}
and $\det DF_K \ge c_0\, h_K^{3}$ a.e. on $\widehat K$.
Moreover, $h\simeq h_K$ for all $K\in\mathcal T_h$, and $\Omega=\bigcup_{K\in\mathcal T_h}K$.

\medskip
\noindent\emph{Parametrizable domains.}
The framework above also applies if $\Omega$ is given as the image of a reference domain.
More precisely, assume there exist a bounded reference domain $\widehat\Omega\subset\R^3$ admitting a
shape-regular simplicial mesh family $\{\widehat{\mathcal T}_h\}_{h\downarrow 0}$ with
$\widehat\Omega=\bigcup_{\widehat K\in\widehat{\mathcal T}_h}\widehat K$, and a bi-Lipschitz map
$\Phi:\widehat\Omega\to\Omega$ with $\Phi\in W^{2,\infty}(\widehat\Omega;\R^3)$ and $\Phi^{-1}\in W^{1,\infty}(\Omega;\R^3)$.
Then, pushing forward the reference mesh by $\Phi$ yields an isoparametric family (cf.\ the standard pushforward construction in \cite{CiarletRaviart1972,Zlamal1973})
$\mathcal T_h:=\{\Phi(\widehat K):\widehat K\in\widehat{\mathcal T}_h\}$ satisfying the above
uniform bounds (possibly with a different constant $\gamma$) and, in particular, $\Omega=\bigcup_{K\in\mathcal T_h}K$. This shows how rich is the family of domains we can deal with.


The (isoparametric) $P_1$ finite element space is defined by pullback: for each $K$, functions in $V_h$
restrict to polynomials of degree $\le 1$ on $\widehat K$ after composition with $F_K^{-1}$.
The operator $I_h:H^2(\Omega)\to V_h$ denotes the nodal Lagrange interpolant (defined using the mesh vertices).












\begin{lemma}[Nodal interpolation estimate for more general domains]
\label{lem:app-interp}
There exists a constant $C_I>0$, independent of $h$, such that for all $w\in H^2(\Omega)$,
\begin{equation}
\label{eq:app-interp}
\|w-I_h w\|_{H^1(\Omega)} \le C_I\, h\,\|w\|_{H^2(\Omega)}.
\end{equation}
\end{lemma}

\begin{proof}
We give a self-contained proof that remains valid for $C^{1,1}$ domains under the mesh assumption above.
Fix $w\in H^2(\Omega)$ and $K\in\mathcal T_h$.
Set $\widehat w:=w\circ F_K\in H^2(\widehat K)$ and denote by $\widehat I$ the nodal $P_1$ interpolant on the reference simplex
$\widehat K$ (with the reference vertices as nodes).
By construction of the isoparametric interpolant, the local interpolant satisfies
\begin{equation}
\label{eq:app-local-interp}
(I_hw)|_{K} = (\widehat I\,\widehat w)\circ F_K^{-1}.
\end{equation}

\smallskip
\noindent\emph{Step 1: reference-element estimate.}
On the fixed simplex $\widehat K$ the Bramble--Hilbert lemma (or an explicit finite-dimensional argument) yields the standard bound
\begin{equation}
\label{eq:app-ref-est}
\|\widehat w-\widehat I\widehat w\|_{H^1(\widehat K)}\le C_{\widehat I}\,|\widehat w|_{H^2(\widehat K)},
\end{equation}
where $C_{\widehat I}$ depends only on the reference simplex.

\smallskip
\noindent\emph{Step 2: mapping of norms.}
Using the change of variables $x=F_K(\widehat x)$, $J_K:=DF_K$, and the chain rule,
\[\nabla_x \big(\phi\circ F_K^{-1}\big)(x)=J_K(\widehat x)^{-\top}\,\nabla_{\widehat x}\phi(\widehat x),\qquad \widehat x=F_K^{-1}(x),\]
we obtain the local $H^1$-estimate
\begin{equation}
\label{eq:app-H1-map}
\|w-I_hw\|_{H^1(K)}\le C\,h_K^{\frac d2-1}\,\|\widehat w-\widehat I\widehat w\|_{H^1(\widehat K)},
\end{equation}
with a constant $C$ depending only on the uniform bounds in \eqref{eq:app-shapereg}.
Similarly, by differentiating twice and using \eqref{eq:app-shapereg}, one has the standard pullback bound
\begin{equation}
\label{eq:app-H2-map}
|\widehat w|_{H^2(\widehat K)}\le C\,h_K^{2-\frac d2}\,\|w\|_{H^2(K)},
\end{equation}
where again $C$ depends only on \eqref{eq:app-shapereg}.
(A direct derivation expands $D^2_{\widehat x}(w\circ F_K)$ into a term involving $D^2_x w$ and a lower-order term involving
$D_x w$ multiplied by $D^2F_K$; the latter is controlled by \eqref{eq:app-shapereg} and absorbed into $\|w\|_{H^2(K)}$.)

\smallskip
\noindent\emph{Step 3: local-to-global.}
Combining \eqref{eq:app-ref-est}--\eqref{eq:app-H2-map} gives
\[
\|w-I_hw\|_{H^1(K)}
\le C\,h_K^{\frac d2-1}\,C_{\widehat I}\,|\widehat w|_{H^2(\widehat K)}
\le C\,C_{\widehat I}\,h_K\,\|w\|_{H^2(K)}.
\]
Summing over $K\in\mathcal T_h$ and using quasi-uniformity ($h_K\simeq h$) yields
\[
\|w-I_hw\|_{H^1(\Omega)}^2
=\sum_{K\in\mathcal T_h}\|w-I_hw\|_{H^1(K)}^2
\le C\,h^2\sum_{K\in\mathcal T_h}\|w\|_{H^2(K)}^2
=C\,h^2\,\|w\|_{H^2(\Omega)}^2,
\]
which implies \eqref{eq:app-interp} with $C_I:=\sqrt{C}$.
The constant depends only on the reference simplex and the uniform $C^1$ shape-regularity bounds \eqref{eq:app-shapereg},
and is therefore independent of $h$.
\end{proof}





\begin{proposition}[Nullspace coercivity on $C^{1,1}$ domains with exact mesh assumption (\ref{eq:app-shapereg})]
\label{prop:app-xi-diverges}
Let $D_M\subset D(\mathcal A)$ be the subspace
\[
D_M:=\{w\in D(\mathcal A):\ \langle d_{M,j},w\rangle = w(x_{h,j})=0\ \forall j=1,\dots,M\}.
\]
Then there exists $c>0$, independent of $h$, such that the constant
\[
\xi_M^+ := \inf\Big\{\frac{\|w\|_{D(\mathcal A)}^2}{\|w\|_V^2}:\ w\in D_M\setminus\{0\}\Big\}
\]
satisfies
\begin{equation}
\label{eq:app-xi-rate}
\xi_M^+ \ge c\,h^{-2}\longrightarrow +\infty\qquad (h\downarrow 0).
\end{equation}
\end{proposition}

\begin{proof}
Let $w\in D_M$. Since $w\in D(\mathcal A)\subset H^2(\Omega)$ and $w(x_{h,j})=0$ for every vertex, the nodal interpolant
satisfies $I_h w\equiv 0$.
Therefore, using Lemma~\ref{lem:app-interp} and Lemma~\ref{lem:app-graphnorm-H2},
\[
\|w\|_{H^1(\Omega)} = \|w-I_h w\|_{H^1(\Omega)} \le C_I h\,\|w\|_{H^2(\Omega)}
\le \frac{C_I}{c_\mathcal A}\, h\,\|w\|_{D(\mathcal A)}.
\]
Since $V=H^1(\Omega)$ with an equivalent norm, we obtain $\|w\|_V\le C h\|w\|_{D(\mathcal A)}$ for a constant $C$ independent
of $h$. Squaring and taking reciprocals yields
\[
\frac{\|w\|_{D(\mathcal A)}^2}{\|w\|_V^2} \ge C^{-2} h^{-2}.
\]
Taking the infimum over $w\in D_M\setminus\{0\}$ gives \eqref{eq:app-xi-rate}.
\end{proof}

\begin{corollary}[Assumption~\ref{assump:KRW-act} on domains satisfying (\ref{eq:app-shapereg})]
\label{cor:app-KRW-act}
With the choices \eqref{eq:app-actuators-dirac} and the biorthogonal family from Lemma~\ref{lem:app-biorthogonal},
Assumption~\ref{assump:KRW-act} holds on bounded $C^{1,1}$ domains, and in particular
$\xi_M^+\to+\infty$.
\end{corollary}

\begin{remark}[Rate in terms of the number of actuators]
If $\{\mathcal T_h\}$ is quasi-uniform, then $M=M(h)\asymp h^{-d}$, hence \eqref{eq:app-xi-rate} implies
$\xi_M^+\gtrsim M^{2/d}$.
\end{remark}





%--------------------------------------------------------
\subsection{Three sufficient mechanisms ensuring solvability of the fixed--point correction}
\label{subsec:fp-solvability-mechanisms}
%--------------------------------------------------------
\noindent
This subsection collects three \emph{complementary} ways to guarantee solvability of the finite-dimensional
fixed--point correction used to compensate the steady bias. The three mechanisms correspond to:
(i) increasing the medium diffusivity $\kappa_m$; (ii) strengthening the stabilizing feedback (gain/scaling $\lambda$);
and (iii) designing the actuator placement to improve conditioning.

%--------------------------------------------------------
\paragraph{Fixed--point equation and a universal contraction criterion.}
%--------------------------------------------------------
Let $\{\varphi_k\}_{k\ge 1}$ be the Neumann eigenfunctions of $-\Delta$ in $\Omega$ with eigenvalues
$0=\lambda_1\le\lambda_2\le\cdots$ and set
\[
X_N:=\mathrm{span}\{\varphi_1,\dots,\varphi_N\},\qquad P_N:H\to X_N \text{ the }H\text{--orthogonal projector},\qquad Q_N:=I-P_N .
\]
In the fixed--point construction of Subsection~\ref{subsec:fixed-point-bias} (notation as there),
the bias operator $T_N:X_N\to X_N$ can be written as
\begin{equation}
\label{eq:TN-abstract}
T_N \;=\; P_N R_{\mathrm{cl}}\, Q_N\, B\,U_N ,
\end{equation}
where:
\begin{itemize}
\item $A y := \kappa_m \Delta y$ (Neumann boundary condition) and $A_{\mathrm{cl}}$ is the closed--loop generator;
\item $S_{\mathrm{cl}}(t)=e^{tA_{\mathrm{cl}}}$ is exponentially stable in $V'$ (see Theorem~\ref{thm:KRW-Vprime});
\item $R_{\mathrm{cl}}:=\displaystyle\int_0^\infty S_{\mathrm{cl}}(t)\,dt$ (so $R_{\mathrm{cl}}=-A_{\mathrm{cl}}^{-1}$ whenever $A_{\mathrm{cl}}$ is invertible);
\item $B:\mathbb{R}^M\to\mathcal{X}$ is the actuator operator (Dirac point actuators, interpreted in the same dual/extrapolation space $\mathcal{X}$ as in Subsection~\ref{subsec:fixed-point-bias});
\item $U_N:X_N\to\mathbb{R}^M$ is the (linear) ``feedforward selector'' used to enforce the cancellation property
\begin{equation}
\label{eq:cancellation-property}
P_N \bigl(A_0 y + B U_N y\bigr)=0,\qquad \forall\,y\in X_N,
\end{equation}
so that the residual forcing generated by the feedforward acts only on the tail.
\end{itemize}
The fixed--point corrected reference $y^\star\in X_N$ is defined by
\begin{equation}
\label{eq:fp-linear-again}
(I+T_N)\,y^\star = y_r \qquad \text{in }X_N.
\end{equation}

\medskip

\noindent
A sufficient condition for unique solvability of \eqref{eq:fp-linear-again} by contraction is $\|T_N\|<1$.
To expose the three mechanisms, we separate the estimate into (i) \emph{low$\leftarrow$high leakage} and
(ii) \emph{conditioning of the feedforward}:
\begin{equation}
\label{eq:TN-factorization}
\|T_N\|_{\mathcal{L}(X_N)}
\;\le\;
\underbrace{\|P_N R_{\mathrm{cl}}Q_N\|_{\mathcal{L}(\mathcal{X})}}_{=: \,L_N}
\;\cdot\;
\|B\|_{\mathcal{L}(\mathbb{R}^M,\mathcal{X})}
\;\cdot\;
\|U_N\|_{\mathcal{L}(X_N,\mathbb{R}^M)}.
\end{equation}

\begin{proposition}[Universal contraction criterion]
\label{prop:universal-contraction}
If
\begin{equation}
\label{eq:universal-contraction}
L_N\,\|B\|\,\|U_N\|<1,
\end{equation}
then $I+T_N$ is invertible on $X_N$, \eqref{eq:fp-linear-again} admits a unique solution $y^\star\in X_N$, and the Picard
iteration $y^{(\ell+1)} := y_r - T_N y^{(\ell)}$ converges geometrically to $y^\star$.
Moreover,
\[
\|y^\star\|_{X_N}\le \frac{1}{1-\|T_N\|}\,\|y_r\|_{X_N}.
\]
\end{proposition}

\begin{proof}
This is immediate from $\|T_N\|<1$ and the Neumann series for $(I+T_N)^{-1}$.
\end{proof}
\noindent
In the remainder of this subsection, we present three independent ways to enforce \eqref{eq:universal-contraction}.
Mechanisms (A) and (B) reduce the leakage factor $L_N$, while mechanism (C) reduces the conditioning factor $\|U_N\|$.

%--------------------------------------------------------
\paragraph{(A) Large diffusivity $\kappa_m$ reduces low$\leftarrow$high leakage.}
%--------------------------------------------------------
We decompose the closed--loop generator with respect to $H=X_N\oplus X_N^\perp$:
\[
A_{\mathrm{cl}}
=
\begin{pmatrix}
A_{11} & A_{12}\\
A_{21} & A_{22}
\end{pmatrix},
\qquad
A_{ij}:=P_i A_{\mathrm{cl}} P_j,\;\; P_1:=P_N,\;\; P_2:=Q_N,
\]
and write
\[
A_{22} = \kappa_m\,Q_N\Delta Q_N + E_{22},\qquad E_{22}:=Q_N\bigl(A_{\mathrm{cl}}-\kappa_m\Delta\bigr)Q_N .
\]
Set $\beta:=\|E_{22}\|_{\mathcal{L}(X_N^\perp)}$ (finite because the feedback/observation terms in $A_{\mathrm{cl}}$ are bounded on the chosen state space).

\begin{lemma}[Tail resolvent estimate]
\label{lem:tail-resolvent-kappa}
If $\kappa_m\lambda_{N+1}>\beta$, then $A_{22}$ is invertible on $X_N^\perp$ and
\begin{equation}
\label{eq:tail-inverse}
\|A_{22}^{-1}\|_{\mathcal{L}(X_N^\perp)}
\;\le\;
\frac{1}{\kappa_m\lambda_{N+1}-\beta}.
\end{equation}
\end{lemma}

\begin{proof}
On $X_N^\perp$, the spectrum of $Q_N\Delta Q_N$ is contained in $\{-\lambda_k\}_{k\ge N+1}$, hence
$\|(\kappa_m Q_N\Delta Q_N)^{-1}\|\le (\kappa_m\lambda_{N+1})^{-1}$.
Write $A_{22}=\kappa_m Q_N\Delta Q_N + E_{22}$ and invert by a Neumann series when we have 
$\|E_{22}(\kappa_m Q_N\Delta Q_N)^{-1}\|<1$, i.e. $\kappa_m\lambda_{N+1}>\beta$.
\end{proof}
\noindent
Assume in addition that the Schur complement on $X_N$,
\[
S:=A_{11}-A_{12}A_{22}^{-1}A_{21},
\]
is invertible and satisfies $\|S^{-1}\|\le C_S$ (typically guaranteed by exponential stability on $X_N$ of the chosen feedback design).
Then the $(1,2)$--block of $A_{\mathrm{cl}}^{-1}$ satisfies the standard identity
\[
P_N A_{\mathrm{cl}}^{-1} Q_N \;=\; -S^{-1}A_{12}A_{22}^{-1}.
\]
Since $R_{\mathrm{cl}}=-A_{\mathrm{cl}}^{-1}$, we obtain:

\begin{proposition}[Leakage bound improved by large $\kappa_m$]
\label{prop:leakage-kappa}
Assume $\kappa_m\lambda_{N+1}>\beta$ and $\|S^{-1}\|\le C_S$. Then
\begin{equation}
\label{eq:leakage-kappa}
L_N=\|P_N R_{\mathrm{cl}}Q_N\|
\;\le\;
\frac{C_S\,\|A_{12}\|}{\kappa_m\lambda_{N+1}-\beta}.
\end{equation}
Consequently, the contraction condition \eqref{eq:universal-contraction} holds whenever
\begin{equation}
\label{eq:kappa-threshold-final}
\frac{C_S\,\|A_{12}\|}{\kappa_m\lambda_{N+1}-\beta}\,\|B\|\,\|U_N\|<1.
\end{equation}
\end{proposition}

\begin{proof}
Combine the block identity with Lemma~\ref{lem:tail-resolvent-kappa} and \eqref{eq:TN-factorization}.
\end{proof}

%--------------------------------------------------------
\paragraph{(B) Stronger stabilizing feedback (gain $\lambda$) reduces $R_{\mathrm{cl}}$ and hence $L_N$.}
%--------------------------------------------------------
A complementary mechanism is to increase the closed--loop decay rate by strengthening the feedback.
Assume the closed--loop semigroup satisfies the exponential stability estimate
\begin{equation}
\label{eq:exp-stab}
\|S_{\mathrm{cl}}(t)\|_{\mathcal{L}(\mathcal{X})} \le M e^{-\alpha t},\qquad t\ge 0,
\end{equation}
for some $M\ge 1$ and $\alpha>0$. Then
\begin{equation}
\label{eq:Rcl-bound}
\|R_{\mathrm{cl}}\|
\;=\;
\Bigl\|\int_0^\infty S_{\mathrm{cl}}(t)\,dt\Bigr\|
\;\le\;
\int_0^\infty \|S_{\mathrm{cl}}(t)\|\,dt
\;\le\;
\frac{M}{\alpha}.
\end{equation}
Since $\|P_N\|,\|Q_N\|\le 1$ on $H$ (and remain bounded in equivalent norms on $X_N$), this implies the crude but robust bound
\begin{equation}
\label{eq:leakage-alpha}
L_N=\|P_NR_{\mathrm{cl}}Q_N\|
\;\le\;
\|R_{\mathrm{cl}}\|
\;\le\;
\frac{M}{\alpha}.
\end{equation}

\begin{proposition}[Contraction criterion in terms of the closed--loop decay rate]
\label{prop:alpha-contraction}
Assume \eqref{eq:exp-stab}. If
\begin{equation}
\label{eq:alpha-contraction}
\frac{M}{\alpha}\,\|B\|\,\|U_N\|<1,
\end{equation}
then the fixed--point problem \eqref{eq:fp-linear-again} has a unique solution and the Picard iteration converges.
\end{proposition}

\begin{proof}
Combine \eqref{eq:leakage-alpha} with Proposition~\ref{prop:universal-contraction}.
\end{proof}

\begin{remark}[How $\lambda$ enters $\alpha$]
In many finite-dimensional designs (pole placement / Riccati design applied to the truncated model on $X_N$),
one can enforce a prescribed decay margin on $X_N$:
\[
\Re\,\sigma\!\bigl(A_{11}(\lambda)\bigr)\le -\alpha_{\mathrm{low}}(\lambda),
\]
where $\alpha_{\mathrm{low}}(\lambda)$ increases when the feedback scaling $\lambda$ increases, provided the pair
$(A_{11},B_{11})$ is stabilizable on $X_N$.
The global $\alpha$ can be taken as $\alpha=\min\{\alpha_{\mathrm{low}}(\lambda),\,\kappa_m\lambda_{N+1}-\beta\}$,
up to a multiplicative constant $M$ (depending on the chosen norm). This makes \eqref{eq:alpha-contraction}
a practical design criterion: increase $\lambda$ until $\alpha$ is large enough.
\end{remark}

%--------------------------------------------------------
\paragraph{(C) Placement design reduces $\|U_N\|$ via conditioning of the sampling matrix.}
%--------------------------------------------------------
The operator $U_N$ is determined by the cancellation requirement \eqref{eq:cancellation-property}.
For point actuators, the relevant finite-dimensional matrix is the sampling (or ``moment'') matrix
\begin{equation}
\label{eq:CN-def}
C_N \;:=\; \bigl(\varphi_k(x_j)\bigr)_{k=1,\dots,N}^{j=1,\dots,M}\in\mathbb{R}^{N\times M},
\end{equation}
where $\{x_j\}_{j=1}^M\subset\Omega$ are the actuator centers. In typical constructions, $U_N$ is chosen as a right-inverse
of $C_N$ (possibly composed with fixed diagonal scalings depending on the exact definition of $B$ and the chosen basis of $X_N$).
The key quantity is the smallest singular value $\sigma_{\min}(C_N)$.

\begin{lemma}[Right-inverse bound]
\label{lem:right-inverse-bound}
Assume $\mathrm{rank}(C_N)=N$ (in particular $M\ge N$). Let $C_N^\dagger$ denote the Moore--Penrose pseudoinverse.
Then $C_N^\dagger$ is a bounded right-inverse, $C_N C_N^\dagger=I_N$, and
\begin{equation}
\label{eq:pseudoinverse-norm}
\|C_N^\dagger\|_{\mathcal{L}(\mathbb{R}^N,\mathbb{R}^M)} = \frac{1}{\sigma_{\min}(C_N)}.
\end{equation}
Consequently, if $U_N$ is chosen proportional to $C_N^\dagger$ (as in the minimum-norm cancellation choice), then
\begin{equation}
\label{eq:UN-bound-smin}
\|U_N\|\;\lesssim\;\frac{1}{\sigma_{\min}(C_N)}.
\end{equation}
\end{lemma}

\begin{proof}
Standard linear algebra: $\|C_N^\dagger\|=1/\sigma_{\min}(C_N)$ for full row rank matrices.
\end{proof}

\begin{proposition}[Contraction via placement conditioning]
\label{prop:placement-contraction}
Assume $\mathrm{rank}(C_N)=N$ and choose $U_N$ according to the minimum-norm cancellation rule so that
\eqref{eq:UN-bound-smin} holds. Then
\begin{equation}
\label{eq:placement-criterion}
\|T_N\| \;\le\; L_N\,\|B\|\,\|U_N\|
\;\lesssim\;
L_N\,\|B\|\,\frac{1}{\sigma_{\min}(C_N)}.
\end{equation}
In particular, a sufficient condition for solvability of \eqref{eq:fp-linear-again} is
\begin{equation}
\label{eq:smin-threshold}
\sigma_{\min}(C_N)\;\gtrsim\; L_N\,\|B\|.
\end{equation}
\end{proposition}

\begin{proof}
Combine \eqref{eq:TN-factorization} with Lemma~\ref{lem:right-inverse-bound}.
\end{proof}

\begin{remark}[Design of placement]
For fixed $N$, the full-row-rank condition $\mathrm{rank}(C_N)=N$ is generic with respect to the point configuration
$\{x_j\}$ (it fails only on a thin set), and oversampling ($M>N$) typically improves $\sigma_{\min}(C_N)$ and robustness.
A practical design principle is to choose $\{x_j\}$ to maximize $\sigma_{\min}(C_N)$ (or, equivalently, minimize $\|C_N^\dagger\|$),
e.g. by greedy selection from a candidate mesh or by D-optimal design maximizing $\det(C_NC_N^\top)$.
\end{remark}

%--------------------------------------------------------
\paragraph{Combined sufficient condition (three options together).}
%--------------------------------------------------------
We summarize the three mechanisms in a single criterion. Combine the $\kappa_m$--driven leakage estimate
\eqref{eq:leakage-kappa} (mechanism (A)) or the decay-rate estimate \eqref{eq:leakage-alpha} (mechanism (B))
with the placement estimate \eqref{eq:UN-bound-smin} (mechanism (C)). For instance, using \eqref{eq:leakage-kappa} and
Lemma~\ref{lem:right-inverse-bound}, a sufficient contraction condition is
\[
\frac{C_S\,\|A_{12}\|}{\kappa_m\lambda_{N+1}-\beta}\;\|B\|\;\frac{1}{\sigma_{\min}(C_N)} \;<\; 1,
\]
which can be enforced by: increasing $\kappa_m$ (stronger tail dissipation), increasing the feedback gain $\lambda$ so that
$C_S$ and/or $\beta$ improve (stronger closed-loop stability), and/or improving the actuator placement to increase $\sigma_{\min}(C_N)$.












\section{Conclusions and perspectives}

\noindent We have combined a discrete effective-medium description of heat generation by
plasmonic nanoparticles with an abstract feedback design for heat-profile
tracking based on point actuators. At the modeling level, we relied on the
rigorous point-approximation framework of \cite{CaoMukherjeeSini2025} to
identify a finite collection of nanoparticles as a system of effective point
actuators for the heat equation: the temperature variation induced by the
particles can be represented, up to small error terms, as a superposition of
heat kernels centered at the nanoparticle locations, with amplitudes solving
a Volterra system driven by the internal electromagnetic energy.

\noindent At the control level, we interpret the heat equation with point actuators in the natural
space $V'$, $V:=D(\mathcal A)$, and consider a feedback law based on the
elliptically smoothed measurements $\mathcal A^{-1}y$.
Using the stabilization theory of \cite{KunRodWal24-cocv}, we obtain exponential decay of the
tracking error in $V'$ under an actuator-family hypothesis, and we recover $H$--tracking after any
positive time $t_0>0$ via parabolic regularization.

\noindent Finally, we coupled the two levels by expressing the effective heat inputs
generated by the nanoparticle ensemble as a linear operator $\mathcal{K}$
acting on a set of electromagnetic illumination parameters. Under suitable
assumptions on the illumination patterns, $\mathcal{K}$ admits a right-inverse
on a relevant finite-dimensional subspace, and the discrepancy between the
ideal feedback inputs and the effective thermo--plasmonic ones can be
quantified in terms of the small geometric parameter $\delta$ and the modeling
error from the discrete effective description. This yields an abstract estimate
on the tracking error between the ideal closed-loop temperature and the
physically realized one, and shows that, in the asymptotic regime considered
in \cite{CaoMukherjeeSini2025}, plasmonic nanoparticles can indeed be used as
actuators for heat-profile tracking.
\bigskip

\noindent Two implementation layers merit further development.
First, while the analysis allows signed coefficients in the abstract intensity parameter, physically one has
$p_\ell(t)=|a_\ell(t)|^2\ge 0$, leading to cone-type feasibility constraints for the reduced actuator input.
Second, the feedback output is formulated via an elliptically regularized point evaluation, which is convenient
for Dirac-actuated heat dynamics but this needs te be linked to more realistic
measurement modalities.
A forthcoming work will incorporate these constraints directly (including nonnegativity and sensor models), and
will provide numerical studies illustrating the closed-loop tracking performance. In addition, we will work directly on the full space avoiding finite-dimensional projections.
\bigskip

\noindent Several directions for further work remain open, including heterogeneous and
structured backgrounds and lower-dimensional configurations with
 robust and optimal designs of nanoparticle configurations
and illuminations, and numerical validation. In addition, we plan to analyze the feedback control part in the whole space, with the diffusion operator $A$ stated now in the whole space with absolutely continuous spectrum. In the same direction, we would like to study tracking other fields, as acoustic fluctuation (in the framework of acoustic cavitation) using either purely acoustic models, see \cite{Mukherjee2024}, in Ultrasound, or multy-physic models as photo-acoustics, see \cite{Ghandriche2022}. Finally, such an approach make also full sense in mechanics where the goal is to generate and track full displacements and not only pressure parts, as in acoustics.    
\bigskip


\noindent
{\bf Acknowledgments:}
The work of S. Rodrigues is funded by national funds through the FCT -- Funda\c{c}\~{a}o para a Ci\^{e}ncia e a Tecnologia, I.P., under the scope of the projects UID/00297/2025 (\url{https://doi.org/10.54499/UID/00297/2025}) and UID/PRR/00297/2025 (\url{https://doi.org/10.54499/UID/PRR/00297/2025}) (Center for Mathematics and Applications -- NOVA Math). The work of M. Sini is partially supported by the Austrian Science Fund (FWF): P 32660 and  P: 36942.  












\bigskip



\begin{thebibliography}{99}

\bibitem{AmmariMillienRuizZhang2017}
H.~Ammari, P.~Millien, M.~Ruiz, and H.~Zhang.
\newblock Mathematical analysis of plasmonic nanoparticles: the scalar case.
\newblock {\em Arch. Ration. Mech. Anal.}, 224(2):597--658, 2017.

\bibitem{AmmariRomeroRuiz2018}
H.~Ammari, F.~Romero, and M.~Ruiz.
\newblock Heat generation with plasmonic nanoparticles.
\newblock {\em Multiscale Model. Simul.}, 16(1):356--384, 2018.

%\bibitem{Armaou2001}
%A.~Armaou and M.~Christofides.
%\newblock Feedback control of {PDE} systems.
%\newblock {\em AIChE Journal}, 42(11):3063--3086, 1996.

\bibitem{AzouaniTiti14}
A. Azouani and E.~S. Titi.
\newblock{Feedback control of nonlinear dissipative systems by finite determining
parameters --- a reaction--diffusion paradigm}. \newblock{\em Evol. Equ. Control Theory}, 3:579--594, 2014.

\bibitem{Barbu2018}
V.~Barbu.
\newblock {\em Controllability and Stabilization of Parabolic Equations}.
\newblock Springer, 2018.

\bibitem{BaffouQuidant2013}
G.~Baffou and R.~Quidant.
\newblock Thermo-plasmonics: using metallic nanostructures as nano-sources of heat.
\newblock {\em Laser Photonics Rev.}, 7(2):171--187, 2013.

\bibitem{BaffouNM2020}
G.~Baffou et~al.
\newblock Applications and challenges of thermoplasmonics.
\newblock {\em Nat. Mater.}, 19:946--958, 2020.

\bibitem{BrennerScottFEM}
S.~C. Brenner and L.~R. Scott,
\emph{The Mathematical Theory of Finite Element Methods},
Texts in Applied Mathematics, Vol.~15, 3rd ed., Springer, 2008.
See Theorem~4.4.20 (and Corollary~4.4.24/Remark~4.4.27) for the $P_1$ nodal
interpolation estimate.


\bibitem{CaoMukherjeeSini2025}
X.~Cao, A.~Mukherjee, and M.~Sini.
\newblock Effective medium theory for heat generation using plasmonics.
\newblock {\em Math. Ann.}, 2025.
\newblock DOI: 10.1007/s00208-025-03300-9.

\bibitem{CasZuazua}
C.~Castro and E.~Zuazua.
\newblock Unique continuation and control for the heat equation from a lower-dimensional moving set.
\newblock {\em SIAM J. Control Optim.}, 43(4):1400--1434, 2005.

\bibitem{Chuang2022}
Y.-C.~Chuang, Y.-T.~Wu, and C.-H.~Chen.
\newblock Recent advances in gold nanomaterials for photothermal cancer therapy.
\newblock {\em Nanomaterials}, 12(2):8, 2022.


\bibitem{CiarletRaviart1972}
P.~G. Ciarlet and P.-A. Raviart.
\newblock Interpolation theory over curved elements, with applications to finite element methods.
\newblock {\em Comput. Methods Appl. Mech. Engrg.}, 1(2):217--249, 1972.


\bibitem{Ghandriche2022}
A. Ghandriche and M. Sini
\newblock Photo-acoustic inversion using plasmonic contrast agents: the full Maxwell model.
\newblock{J. Differential Equations} 341 (2022), 1–78.


\bibitem{EggerFritzKunRod25}
H. Egger and M. Fritz and K. Kunisch and S.~S. Rodrigues.
\newblock{On feedback stabilization for the Cahn–Hilliard equation and its numerical
approximation}.
\newblock{\em IMA J. Numer. Anal.} (to appear), 2025. \url{https://doi.org/10.1093/imanum/draf118}


%\bibitem{EngelNagel-Book}
%K.-J.~Engel and R.~Nagel.
%\newblock {\em One-Parameter Semigroups for Linear Evolution Equations}.
%\newblock Graduate Texts in Mathematics, Vol.~194, Springer-Verlag, New York, 2000.

\bibitem{FardigolaKhalina2022}
L.~Fardigola and K.~Khalina.
\newblock Controllability problems for the heat equation in a half-plane controlled by the Dirichlet boundary condition with a point-wise control.
\newblock {\em J. Math. Phys. Anal. Geom.}, 18(1):75--104, 2022.

\bibitem{Hamidoglu2013}
A.~Hamidoglu.
\newblock Null controllability of the heat equation with switching pointwise controls.
\newblock {\em Trans. Inst. Math. Mech.}, 35(4):499--512, 2013.

\bibitem{Jorgensen2016}
J.~T. J{\o}rgensen et~al.
\newblock Single particle and {PET}-based platform for identifying optimal plasmonic nano-heaters for photothermal cancer therapy.
\newblock {\em Sci. Rep.}, 6:30076, 2016.

\bibitem{Kim2023}
D.~Kim et~al.
\newblock Effect of gold nanoparticle distribution radius on photothermal therapy of skin cancer: a numerical study.
\newblock {\em Sci. Rep.}, 13:39040, 2023.

\bibitem{KunRodWal24-cocv}
K.~Kunisch, S.~S.~Rodrigues, and D.~Walter.
\newblock Stabilizability for nonautonomous linear parabolic equations with actuators as distributions.
\newblock {\em ESAIM: Control Optim. Calc. Var.}, 30:43, 2024.

\bibitem{KrsticSmyshlyaev2008}
M.~Krstic and A.~Smyshlyaev.
\newblock {\em Boundary Control of PDEs: A Course on Backstepping Designs}.
\newblock SIAM, 2008.

\bibitem{Lenoir1986}
M.~Lenoir.
\newblock Optimal isoparametric finite elements and error estimates for domains involving curved boundaries.
\newblock {\em SIAM J. Numer. Anal.}, 23(3):562--580, 1986.

\bibitem{Letrouit2019}
C.~Letrouit.
\newblock From internal to pointwise control for the 1{D} heat equation and minimal control time.
\newblock {\em Systems Control Lett.}, 133:104549, 2019.

%\bibitem{LionsMagenesI}
%J.-L.~Lions and E.~Magenes.
%\newblock {\em Non-Homogeneous Boundary Value Problems and Applications. Vol.~I}.
%\newblock Die Grundlehren der Mathematischen Wissenschaften, Vol.~181, Springer-Verlag, 1972.

\bibitem{Lunardi-Book}
A.~Lunardi.
\newblock {\em Analytic Semigroups and Optimal Regularity in Parabolic Problems}.
\newblock Progress in Nonlinear Differential Equations and Their Applications, Vol.~16, Birkh\"auser, Basel, 1995.

\bibitem{LunasinTiti17}
E. Lunasin and E.~S. Titi.
\newblock{Finite determining parameters feedback control for distributed nonlinear dissipative
systems—a computational study}. \newblock{\em Evol. Equ. Control Theory}, 6:535--557, 2017.

\bibitem{Martinez2015}
J.~de~Jes{\'u}s Mart{\'i}nez.
\newblock {\em Modeling and Controllability of a Heat Equation with a Point Control}.
\newblock PhD thesis, Iowa State University, 2015.

\bibitem{Mukherjee2022}
A.~Mukherjee and M.~Sini.
\newblock Heat generation using Lorentzian nanoparticles. The full Maxwell system.
\newblock {\em SIAM J. Appl. Math.}, 84(1):285--315, 2024.

\bibitem{Mukherjee2024}
A.~Mukherjee and M.~Sini.
\newblock Time-dependent acoustic waves generated by multiple resonant bubbles: application to acoustic cavitation.
\newblock{J. Evol. Equ.} 24 (2024), no. 4, Paper No. 90, 40 pp.

\bibitem{Pazy1983}
A.~Pazy.
\newblock {\em Semigroups of Linear Operators and Applications to Partial Differential Equations}.
\newblock Springer, 1983.

\bibitem{PhotothermalJAP2016}
S.~J. Norton and T.~Vo-Dinh.
\newblock Photothermal effects of plasmonic metal nanoparticles in a fluid: a modeling study.
\newblock {\em J. Appl. Phys.}, 119(8):083105, 2016.

\bibitem{RaymondBook}
J.-P. Raymond.
\newblock {\em Optimal Control of Partial Differential Equations}.
\newblock Lecture notes, 2016.

\bibitem{Rodrigues2021}
S.~S. Rodrigues.
\newblock Oblique projection output-based feedback exponential stabilization of nonautonomous parabolic equations.
\newblock {\em Automatica}, 129:109621, 2021.

\bibitem{TucsnakWeiss2009}
M.~Tucsnak and G.~Weiss.
\newblock {\em Observation and Control for Operator Semigroups}.
\newblock Birkh\"auser, 2009.

\bibitem{Vines2019}
J.~B. Vines et~al.
\newblock Gold nanoparticles for photothermal cancer therapy.
\newblock {\em Front. Pharmacol.}, 10:1673, 2019.

\bibitem{WuSheldon2023}
S.~Wu and M.~Sheldon.
\newblock Mechanisms of photothermalization in plasmonic nanostructures: insights into the steady state.
\newblock {\em Annu. Rev. Phys. Chem.}, 74:195--216, 2023.

\bibitem{StrangBanks}
G.~Strang and T.~Nguyen.
\newblock {\em Wavelets and Filter Banks}.
\newblock Wellesley-Cambridge Press, 1996.

%\bibitem{Temam-Book}
%R.~Temam.
%\newblock {\em Infinite-Dimensional Dynamical Systems in Mechanics and Physics}.
%\newblock Applied Mathematical Sciences, Vol.~68, Springer-Verlag, New York, 1988.




\bibitem{Zlamal1973}
M.~Zlamal.
\newblock Curved elements in the finite element method. {I}.
\newblock {\em SIAM J. Numer. Anal.}, 10(1):229--240, 1973.


\end{thebibliography}

\end{document}











\medskip
\noindent\textbf{Divergence of the nullspace coercivity constant.}
We take $\Omega\subset\R^3$ to be a general bounded \textcolor{blue}{piecewise~$C^{2}$ domain in a class generalizing polyhedral domains (and including~\noteC{I would believe this true for\eqref{eq:app-shapereg}! do we know some ref/example?} (some/all?)~$C^2$ domains), by following a variation of arguments in~\cite{AzouaniTiti14,LunasinTiti17,EggerFritzKunRod25} and by working} with a boundary-fitted isoparametric mesh family
$\{\mathcal T_h\}_{h\downarrow 0}$: each element $K\in\mathcal T_h$ is the image of a fixed reference simplex
$\widehat K$ under a bijective map $F_K\in W^{2,\infty}(\widehat K;\R^3)$.
We assume uniform (isoparametric) shape-regularity and quasi-uniformity: there exist constants $\gamma,c_0>0$,
independent of $h$ and $K$, such that
\begin{equation}\label{eq:app-shapereg}
\|DF_K\|_{L^\infty(\widehat K)}\le \gamma h_K,\qquad
\|(DF_K)^{-1}\|_{L^\infty(\widehat K)}\le \gamma h_K^{-1},\qquad
\|D^2F_K\|_{L^\infty(\widehat K)}\le \gamma h_K,
\end{equation}
and $\det DF_K \ge c_0\, h_K^{3}$ a.e. on $\widehat K$.
Moreover, $h\simeq h_K$ for all $K\in\mathcal T_h$, and $\Omega=\bigcup_{K\in\mathcal T_h}K$.

The (isoparametric) $P_1$ finite element space is defined by pullback: for each $K$, functions in $V_h$
restrict to polynomials of degree $\le 1$ on $\widehat K$ after composition with $F_K^{-1}$.
The operator $I_h:H^2(\Omega)\to V_h$ denotes the nodal Lagrange interpolant (defined using the mesh vertices).
































%========================================================



\section{How Theorem~\ref{thm:5.7} answers the original control problem}
\label{rem:link-original-problem}
We divide it into few steps:

\begin{enumerate}
\item The original control question can be formulated as follows: given a desired
temperature profile $y_r$ in $\Omega$, can one design incident electromagnetic
fields acting on a finite collection of thermo--plasmonic nanoparticles so that
the resulting temperature $y_{\mathrm{phys}}(t)$ in the medium tracks $y_r$?

Our analysis proceeds in two steps.

\smallskip
\noindent
(1) \emph{Abstract heat tracking with point actuators.}
In Section~\ref{sec:stab-KRW} we use the feedback law, augmented by a constant feedforward
$u_r$ for non-constant references (see \eqref{eq:ref-finite-modes}--\eqref{eq:feedforward-linear-system}).
The resulting closed loop for the error $z(t)=y(t)-y_r$ is the affine system \eqref{eq:KRW-closed-loop}.
Theorem~\ref{thm:KRW-Vprime} yields exponential stabilization in the $V'$ framework of the shifted error
$z(t)-z_\infty$, and hence convergence of $y_{\mathrm{ideal}}(t)$ towards $y_r+z_\infty$ in $V'$ (and, by smoothing,
in $H$ for any $t\ge t_0>0$).

\smallskip
\noindent
(2) \emph{Realization of the ideal actuators by thermo--plasmonic nanoparticles.}
Sections~\ref{sec:model} and \ref{sec:coupling} show that the cluster of plasmonic nanoparticles, illuminated
by incident fields $p$, generates an effective heat input $G=\mathcal{K}p+\mathcal{R}^\delta p$ in
the background heat equation. Proposition~\ref{prop:K-right-inverse} ensures that on a finite
dimensional subspace $Y_0$ we can invert the leading part of $\mathcal{K}$: for any
desired abstract input $u_{\mathrm{des}}\in Y_0$ there exists an incident field
$p=\mathcal{K}^\dagger u_{\mathrm{des}}$ such that $G=u_{\mathrm{des}}+\mathcal{R}^\delta p$ with
$\|\mathcal{R}^\delta p\|_{L^2(0,T;\mathbb{R}^M)}\le C_T\delta^\mu \|u_{\mathrm{des}}\|$.

\smallskip
Theorem~\ref{thm:5.7} combines these two steps. It compares the ``ideal''
closed-loop temperature $y_{\mathrm{ideal}}$ driven by the abstract control
$u_{\mathrm{des}}$ with the ``physical'' temperature $y_{\mathrm{phys}}$
driven by the thermo--plasmonic actuators corresponding to
$p=\mathcal{K}^\dagger u_{\mathrm{des}}$. The estimate
\[
\sup_{t\in[0,T]} \|y_{\mathrm{phys}}(t)-y_{\mathrm{ideal}}(t)\|_H
  \le C_T\,\eta(\delta)\,\|u_{\mathrm{des}}\|_{L^2(0,T;\mathbb{R}^M)},
\]
with $\eta(\delta)\to 0$ as $\delta\to 0$, shows that the physical temperature
follows the ideal one up to an error that can be made arbitrarily small by
choosing sufficiently small particles (i.e.\ working within the validity
regime of the effective model). Since the closed loop yields exponential decay of the shifted error $z_{\mathrm{ideal}}(t)-z_\infty$ in $V'$ (Theorem~\ref{thm:KRW-Vprime}), combining Theorem~\ref{thm:5.7} and Corollary~\ref{cor:total-error} shows that the nanoparticle--based system drives $y_{\mathrm{phys}}(t)$ towards $y_r+z_\infty$ up to a discrepancy that can be made small by choosing $Y_0$ rich enough and by taking $\delta$ sufficiently small.

In this sense, Theorem~\ref{thm:5.7} provides a rigorous and quantitative
affirmative answer to the original problem of using thermo--plasmonic
nanoparticles as actuators for heat tracking in $\Omega$.

\medskip

\item We now explain how the restriction $u_{\mathrm{des}}\in Y_0$ can be enforced
for a given reference profile $y_r$. Let $u_{\mathrm{ideal}}$ denote the ideal Dirac-actuator input generated by the feedback
(Section~\ref{sec:stab-KRW}), i.e.
\[
  u_{\mathrm{ideal},j}(t)
  := (u_r)_j-\lambda\,\big(\mathcal A^{-1}(y_{\mathrm{ideal}}(t)-y_r)\big)(x_j),
  \qquad j=1,\dots,M,
\]
where $y_{\mathrm{ideal}}$ is the corresponding closed-loop trajectory.
In general $u_{\mathrm{ideal}}$ lies in $L^2(0,T;\mathbb{R}^M)$, not necessarily in the
finite-dimensional subspace $Y_0$ that can be realized efficiently by the thermo--plasmonic actuators.

A natural way to enforce the constraint $u_{\mathrm{des}}\in Y_0$ is to
project $u_{\mathrm{ideal}}$ onto $Y_0$ and to use only this projected signal
as the ``desired'' input for the thermo--plasmonic realization. More
precisely, we set
\begin{equation}
\label{eq:5.Y0-u-des}
  u_{\mathrm{des}} := P_{Y_0} u_{\mathrm{ideal}},
\end{equation}
where $P_{Y_0}$ denotes the orthogonal projection from
$L^2(0,T;\mathbb{R}^M)$ onto $Y_0$. By construction, $u_{\mathrm{des}}\in Y_0$
and therefore satisfies the hypothesis of Theorem~\ref{thm:5.7}.

Let $y_{Y_0}$ denote the solution of the heat equation driven by
$u_{\mathrm{des}}$:
\begin{equation}
\label{eq:5.Y0}
\begin{cases}
\dot y_{Y_0}(t) = A y_{Y_0}(t) + B u_{\mathrm{des}}(t), & t>0,\\[0.2em]
y_{Y_0}(0) = y_0.
\end{cases}
\end{equation}
The next proposition quantifies the discrepancy between $y_{Y_0}$ and the
ideal closed-loop trajectory $y_{\mathrm{ideal}}$.

\begin{proposition}[Effect of restricting the control to $Y_0$ (in the $V'$ framework)]
\label{prop:projection-error}
Let $T>0$ and let $u_{\mathrm{ideal}}\in L^2(0,T;\R^M)$. Define
$u_{\mathrm{des}}:=P_{Y_0}u_{\mathrm{ideal}}$ as in \eqref{eq:5.Y0-u-des}.
Let $y_{\mathrm{ideal}}$ be the solution of \eqref{eq:heat-DAdual} driven by $u_{\mathrm{ideal}}$
and let $y_{Y_0}$ be the solution of \eqref{eq:5.Y0} driven by $u_{\mathrm{des}}$, with the same initial
condition $y_0\in H\subset V'$. Then there exists $C_T>0$ such that
\begin{equation}
\label{eq:proj-error}
\| y_{Y_0} - y_{\mathrm{ideal}} \|_{C([0,T];V')}^2
+ \| y_{Y_0} - y_{\mathrm{ideal}} \|_{L^2(0,T;H)}^2
\le C_T\,\|(I-P_{Y_0})u_{\mathrm{ideal}}\|_{L^2(0,T;\R^M)}^2.
\end{equation}
Moreover, for every $t_0\in(0,T]$ there exists $C_{T,t_0}>0$ such that
\begin{equation}
\label{eq:proj-error-H-after-t0}
\sup_{t\in[t_0,T]} \| y_{Y_0}(t) - y_{\mathrm{ideal}}(t) \|_{H}
\le C_{T,t_0}\,\|(I-P_{Y_0})u_{\mathrm{ideal}}\|_{L^2(0,T;\R^M)}.
\end{equation}
In particular,
\begin{equation}
\label{eq:best-approx}
\| y_{Y_0} - y_{\mathrm{ideal}} \|_{C([0,T];V')}
\le C_T^{1/2}\inf_{v\in Y_0}\|u_{\mathrm{ideal}}-v\|_{L^2(0,T;\R^M)}.
\end{equation}
\end{proposition}

\begin{proof}
Set $e_{\mathrm{proj}}:=y_{Y_0}-y_{\mathrm{ideal}}$. Then
\[
\dot e_{\mathrm{proj}}(t) = A e_{\mathrm{proj}}(t) + B\bigl(u_{\mathrm{des}}(t)-u_{\mathrm{ideal}}(t)\bigr)
= A e_{\mathrm{proj}}(t) + B\,(I-P_{Y_0})u_{\mathrm{ideal}}(t),
\qquad e_{\mathrm{proj}}(0)=0,
\]
in $D(\mathcal A)'$. Applying Proposition~\ref{prop:wp-Dirac} yields \eqref{eq:proj-error}.
The estimate \eqref{eq:proj-error-H-after-t0} follows from \eqref{eq:proj-error} and
Corollary~\ref{cor:L2-after-t0}. Finally, since $P_{Y_0}$ is the orthogonal projector onto $Y_0$,
\[
\|(I-P_{Y_0})u_{\mathrm{ideal}}\|_{L^2} = \inf_{v\in Y_0}\|u_{\mathrm{ideal}}-v\|_{L^2},
\]
which gives \eqref{eq:best-approx}.
\end{proof}

\medskip
\item Theorem~\ref{thm:5.7} then compares $y_{Y_0}$ with the physical temperature
$y_{\mathrm{phys}}$ generated by the thermo--plasmonic actuators when the
electromagnetic control is chosen as $p = \mathcal{K}^\dagger u_{\mathrm{des}}$.
Combining Proposition~\ref{prop:projection-error} with Theorem~\ref{thm:5.7} yields an
explicit bound on the total tracking error between the ideal closed-loop
trajectory and the physically realized one.

\begin{corollary}[Total tracking error: projection + nanoparticles (in $V'$)]
\label{cor:total-error}
Let $T>0$ and let $u_{\mathrm{ideal}}\in L^2(0,T;\R^M)$ be an ``ideal'' heat input.
Set $u_{\mathrm{des}}:=P_{Y_0}u_{\mathrm{ideal}}$ and let $y_{\mathrm{ideal}}$, $y_{Y_0}$ be as in
Proposition~\ref{prop:projection-error}. Let $y_{\mathrm{phys}}$ be the physical temperature produced by
the thermo--plasmonic actuators when the intensity control $p$ is chosen as in Theorem~\ref{thm:5.7}:
\[
  p = \mathcal{K}^\dagger u_{\mathrm{des}}.
\]
Then there exists $C_T>0$ such that
\begin{equation}
\label{eq:total-error}
\|y_{\mathrm{phys}}-y_{\mathrm{ideal}}\|_{C([0,T];V')}^2 + \|y_{\mathrm{phys}}-y_{\mathrm{ideal}}\|_{L^2(0,T;H)}^2
\le C_T\Bigl(
\|(I-P_{Y_0})u_{\mathrm{ideal}}\|_{L^2(0,T;\R^M)}^2
+ \eta(\delta)^2\,\|u_{\mathrm{des}}\|_{L^2(0,T;\R^M)}^2
\Bigr),
\end{equation}
where $\eta(\delta)\to 0$ is as in Theorem~\ref{thm:5.7}. Moreover, for every $t_0\in(0,T]$,
\[
\sup_{t\in[t_0,T]}\|y_{\mathrm{phys}}(t)-y_{\mathrm{ideal}}(t)\|_{H}
\le C_{T,t_0}\Bigl(
\|(I-P_{Y_0})u_{\mathrm{ideal}}\|_{L^2(0,T;\R^M)}
+ \eta(\delta)\,\|u_{\mathrm{des}}\|_{L^2(0,T;\R^M)}
\Bigr),
\]
for some constant $C_{T,t_0}>0$.
\end{corollary}

\begin{proof}
We decompose
\[
y_{\mathrm{phys}}-y_{\mathrm{ideal}}=(y_{\mathrm{phys}}-y_{Y_0})+(y_{Y_0}-y_{\mathrm{ideal}}).
\]
Applying Theorem~\ref{thm:5.7} to $y_{\mathrm{phys}}-y_{Y_0}$ and Proposition~\ref{prop:projection-error}
to $y_{Y_0}-y_{\mathrm{ideal}}$, and then using the triangle inequality, yields \eqref{eq:total-error}.
The $H$--estimate on $[t_0,T]$ follows similarly from \eqref{eq:tracking-estimate-H-after-t0} and
\eqref{eq:proj-error-H-after-t0}.
\end{proof}
\end{enumerate}

\medskip

Proposition~\ref{prop:projection-error} measures
the error introduced by restricting the abstract feedback design to controls
lying in the finite-dimensional subspace $Y_0$, while Theorem~\ref{thm:5.7} measures
the additional error incurred when realizing such controls by thermo--plasmonic
nanoparticles. Corollary~\ref{cor:total-error} shows that the total tracking
error is the sum of these two contributions: a purely ``approximation'' term,
given by the best $L^2$--approximation of $u_{\mathrm{ideal}}$ in $Y_0$, and
a ``physical realization'' term of order $\eta(\delta)$ coming from the
effective thermo--plasmonic model.






%\paragraph{Implementation of the intensity control $p$.}
%In physical terms $p_\ell(t)=|a_\ell(t)|^2\ge 0$. The abstract framework below treats $p$ as an element of a linear
%space (hence allows signed components). This can be reconciled with nonnegativity by either restricting the desired
%effective inputs to those attainable with nonnegative $p$, or (more generally) by introducing two nonnegative intensity
%channels $p=p^{+}-p^{-}$ (componentwise), realized by time-multiplexing two copies of the illumination dictionary.
%In the remainder of the paper, $p$ denotes this signed intensity parameter.
